%=========================================================================
\chapter{Code Analysis}
\label{chap:analysis}
\renewcommand{\sectiontitle}{}
%=========================================================================

This chapter briefly describes all static code analyses that are integral part of the \gls{WIR} library.
\cref{sec:containers} first introduces all generic base classes that form the analysis framework within \gls{WIR}.
Hereafter, \crefrange{sec:cflow}{sec:flowfacts} are devoted to the control flow analyses, while \crefrange{sec:reachdef}{sec:bitdflow} present all data flow analyses.


%-------------------------------------------------------------------------
\section{Generic Analysis Framework}
\label{sec:containers}
%-------------------------------------------------------------------------

All code analyses presented in the rest of this chapter inherit from the base class \texttt{WIR\_Analysis} declared in file \texttt{WIR/analyses/generic/wiranalysis.h} so that all actual analyses share the very same public \gls{API} and thus are easy to use.
This public interface is shown in \cref{lst:wiranalysis}.

\begin{lstlisting}[style=code,float=ht,caption={Code Analysis API of Class \texttt{WIR\_Analysis}},label=lst:wiranalysis]
// Default constructor for system-level analysis.
explicit WIR_Analysis::WIR_Analysis@(@ WIR_System æ&æs @)@;Ü\label{lst:wiranalysis1}Ü

// Default constructor for compilation unit-level analysis.
explicit WIR_Analysis::WIR_Analysis@(@ WIR_CompilationUnit æ&æc @)@;Ü\label{lst:wiranalysis2}Ü

// Default constructor for function-level analysis.
explicit WIR_Analysis::WIR_Analysis@(@ WIR_Function æ&æf @)@;Ü\label{lst:wiranalysis3}Ü

// Default constructor for basic block-level analysis.
explicit WIR_Analysis::WIR_Analysis@(@ WIR_BasicBlock æ&æb @)@;Ü\label{lst:wiranalysis4}Ü

// Default constructor for instruction-level analysis.
explicit WIR_Analysis::WIR_Analysis@(@ WIR_Instruction æ&æi @)@;Ü\label{lst:wiranalysis5}Ü

// Default constructor for operation-level analysis.
explicit WIR_Analysis::WIR_Analysis@(@ WIR_Operation æ&æo @)@;Ü\label{lst:wiranalysis6}Ü

// Performs the actual code analysis.
@void@ WIR_Analysis::analyze@( void )@;Ü\label{lst:wiranalysis7}Ü
\end{lstlisting}

As can be seen, various constructors can be chosen which determine whether a code analysis shall be carried out for a complete \gls{WIR} system (cf. \cref{lst:wiranalysis1}), or just for a single compilation unit, function, basic block etc. (\crefrange{lst:wiranalysis2}{lst:wiranalysis6}).
After an analysis has been set up this way, only the \texttt{analyze} method from \cref{lst:wiranalysis7} needs to be invoked in order to apply the actual analysis to the \gls{WIR} object used during construction.

\cref{lst:wiranalysis} furthermore shows that there are no methods at all that allow to get and to inspect analysis data that has been produced by the \texttt{analyze} method.
This in turn is motivated by the fact that any kind of meta-data that provides additional information about a system's code properties is modeled within \gls{WIR} by so-called containers.
Consequently, all code analyses presented in the following produce such containers and attach them to the various classes of the \gls{WIR} software model from \cref{fig:swmodel}, ranging from a top-level system down to the various flavors of \gls{WIR} parameters.
\cref{tab:anacont} briefly summarizes all \gls{WIR} analyses and provides an overlook over the actual container classes that are generated by the respective analyses as well as over those \gls{WIR} classes to which these containers get attached.
The following sections will provide in-depth information about the meaning of all these containers and how to query them for their analysis data in detail.

\begin{wirNote}
  Whenever containers are attached to \gls{WIR} objects by an anysis' \texttt{analyze} method, the analysis data stored in these containers reflects the state of the control and data flow at exactly that point when the analysis was carried out.
  If \gls{WIR} objects are modified after an invocation of \texttt{analyze}, potentially existing containers are not automatically updated so that the state of \gls{WIR} and that of the attached containers can be incoherent.
  Application code thus either explicitly has to re-run an analysis from scratch, or the data inside attached containers needs to be updated incrementally after some \gls{WIR} modification.
\end{wirNote}

\begin{table}[t]
  \begin{center}
    \caption{Analysis Containers and their Fixing Points Within WIR}
    \label{tab:anacont}
    \begin{tabular}{|ll|l|l|l|}
      \hline
      & \textbf{Analysis} & \textbf{Container Class} & \textbf{Fixing Points} \\
      \hhline{|==|=|=|}
      \parbox[t]{4mm}{\multirow{5}{*}{\rotatebox[origin=c]{90}{\textbf{Control Flow}}}} & Reachability & \texttt{WIR\_Reachability} & Basic Blocks \\
      & Domination & \texttt{WIR\_Domination} & Basic Blocks and Instructions \\
      & Immediate Domination & \texttt{WIR\_ImmediateDomination} & Basic Blocks \\
      & Structural Control Flow & \texttt{WIR\_ControlTree} & Basic Blocks \\
      & Flow Facts & \texttt{WIR\_FlowFactRef} & Basic Blocks \\
      \hhline{|==|=|=|}
      \parbox[t]{4mm}{\multirow{5}{*}{\rotatebox[origin=c]{90}{\textbf{Data Flow}}}} & Reaching Definitions & \texttt{WIR\_ReachingDefinitions} & Instructions \\
      & Available Definitions & \texttt{WIR\_AvailableDefinitions} & Instructions \\
      & Def-Use/Use-Def Chains & \texttt{WIR\_DUUDChain} & Register Parameters \\
      & Liveness & \texttt{WIR\_LiveOut} & Instructions \\
      & Bit-True Data Flow & \texttt{WIR\_BitValues} & Register and Immediate Parameters \\
      \hline
    \end{tabular}
  \end{center}
\end{table}

All \gls{WIR} classes acting as ``fixing points'' for such containers share a common \gls{API} in order to generically attach or access containers.
This is achieved by having all these classes inherit from the base class \texttt{WIR\_Container\_API} as depicted in \cref{fig:containerapi}.
This \gls{API} is declared in file \texttt{WIR/wir/API/wircontainerapi.h}, and it provides the methods shown in \cref{lst:contapiins} for the insertion and destruction of containers.

\begin{figure}[t]
  \begin{center}
    \includegraphics[width=0.91\textwidth]{containerAPI.png}
  \end{center}
  \caption{Inheritance Diagram of Class \texttt{WIR\_Container\_API}}
  \label{fig:containerapi}
\end{figure}

\begin{lstlisting}[style=code,float=!b,caption={Insertion and Removal of WIR Containers},label=lst:contapiins]
// Copy-insert a new container.
@void@ WIR_Container_API::insertContainer@(@ const WIR_BaseContainer æ&æc @)@;Ü\label{lst:contapiins1}Ü

// Move-insert a new container.
@void@ WIR_Container_API::insertContainer@(@ WIR_BaseContainer æ&&æc @)@;Ü\label{lst:contapiins2}Ü

// Insert a new container.
@void@ WIR_Container_API::insertContainer@(@ WIR_BaseContainer æ*æp @)@;Ü\label{lst:contapiins3}Ü

// Remove all containers of the specified container type (recursively).
@void@ WIR_Container_API::eraseContainers@(@ WIR_id_t id, @bool@ r æ=æ false @)@;Ü\label{lst:contapiins4}Ü

// Remove the specified container.
@void@ WIR_Container_API::eraseContainer@(@ const WIR_BaseContainer æ&æc @)@;Ü\label{lst:contapiins5}Ü

// Remove all containers.
@void@ WIR_Container_API::clearContainers@( void )@;Ü\label{lst:contapiins6}Ü
\end{lstlisting}

The \texttt{insert} methods from \crefrange{lst:contapiins1}{lst:contapiins3} allow to attach a new container to some \gls{WIR} object.
Similar to \cref{chap:lists}, there are dedicated methods for the copy- and move-insertion of containers.
Furthermore, one of the very rare cases within \gls{WIR} where a pointer shows up in the public \gls{API} can be seen in \cref{lst:contapiins3}.
Here, the user is responsible to dynamically create a new container and to pass a pointer to this container to \texttt{insertContainer}.

\begin{wirWarn}
  After such a pointer-based insertion, class \texttt{WIR\_Container\_API} takes over full control over the ownership of that pointer.
  In particular, it automatically destroys the object pointed to during erasure.
  Thus, the pointer passed to \texttt{insertContainer} should not be used by application code any more afterwards.
\end{wirWarn}

Method \texttt{eraseContainers} (cf. \cref{lst:contapiins4}) can be used to remove all containers of a given container class shown in \cref{tab:anacont}.
For example, \texttt{f.eraseContainers( WIR\_Domination::getContainerTypeID(), true )} would have the effect that all containers holding domination information will be removed from \gls{WIR} function \texttt{f}.
The Boolean flag serving as second argument of this method has the effect that the code hierarchy underneath, e.g., function \texttt{f} will also be inspected and containers be removed.
Method \texttt{eraseContainer} (\cref{lst:contapiins5}) can be used to remove one given container, and \texttt{clearContainers} (\cref{lst:contapiins6}) removes all containers attached to some \gls{WIR} object.
All these \texttt{erase} methods automatically delete the corresponding container objects from memory.

All methods to look up containers attached to some \gls{WIR} object are summarized in \cref{lst:contapilookup}.
\texttt{getContainer\-Types} from \cref{lst:contapilookup1} provides the set of unique numerical IDs of all different container classes that are currently attached.
The full set of all currently attached containers can be retrieved using method \texttt{getContainers} (\cref{lst:contapilookup2}), its return type \texttt{WIR\_ContainerSet} is declared as \texttt{std::set<std::reference\_wrapper<WIR\_BaseContainer>, WIR\_Compare<WIR\_BaseContainer>{}>} (see \cref{chap:setsmaps}).
In analogy to method \texttt{WIR\_System::getComponents} (cf. \cref{lst:syshw2} of \cref{lst:syshw}), \texttt{getContainers} from \cref{lst:contapilookup3} provides a set of (wrapped) references to only those containers that can properly be casted to the specified type name \texttt{T}.
The different \texttt{contains} methods can eventually be used to check whether containers are attached to some \gls{WIR} object at all (\cref{lst:contapilookup4}), to check whether containers of a given class from \cref{tab:anacont} are attached (\cref{lst:contapilookup5}), or whether a specific container is attached (\cref{lst:contapilookup6}).

\begin{lstlisting}[style=code,float=t,caption={Lookup of WIR Containers},label=lst:contapilookup]
// Determine all types of containers currently attached to some WIR object.
ästd::setäæ<æWIR_id_tæ>æ WIR_Container_API::getContainerTypes@( void )@ const;Ü\label{lst:contapilookup1}Ü

// Get the complete set of references to containers.
const WIR_ContainerSet WIR_Container_API::getContainers@( void )@ const;Ü\label{lst:contapilookup2}Ü

// Get the set of all containers of a specified type.
templateæ<ætypename Tæ>æ
const ästd::setäæ<æästd::reference_wrapperäæ<æTæ>æ, WIR_Compareæ<æTæ>>æ WIR_Container_API::getContainers@( void )@ const;Ü\label{lst:contapilookup3}Ü

// Determine whether some containers are attached at all.
@bool@ WIR_Container_API::containsContainers@( void )@ const;Ü\label{lst:contapilookup4}Ü

// Determine whether containers of a specified type are attached.
@bool@ WIR_Container_API::containsContainers@(@ WIR_id_t id @)@ const;Ü\label{lst:contapilookup5}Ü

// Determine whether a given container is attached.
@bool@ WIR_Container_API::containsContainer@(@ const WIR_BaseContainer æ&æc @)@ const;Ü\label{lst:contapilookup6}Ü
\end{lstlisting}

\cref{lst:contapiins,lst:contapilookup} show that generic containers of the abstract base class \texttt{WIR\_BaseContainer} are used.
Thus, all analysis containers from \cref{tab:anacont} are derived from this base class.
\cref{lst:basecontainer} shows the relevant parts of the public interface of this base class.
As \texttt{WIR\_BaseContainer} contains purely virtual methods, no actual objects of this class can be created.
We thus omit the description of constructors here for the sake of brevity.

Some of the previously introduced methods like, e.g., \texttt{eraseContainers} (cf. \cref{lst:contapiins4} of \cref{lst:contapiins}) or \texttt{getCon\-tainerTypes} (\cref{lst:contapilookup1} of \cref{lst:contapilookup})  dealt with unique numerical IDs of the types of containers.
These container type IDs have the effect that, e.g., \texttt{WIR\_ControlTree} containers can generally be distinguished from \texttt{WIR\_LiveOut} containers without having to do dynamic type casts in C++.
The two variants of the \texttt{getContainer\-Type} method from \cref{lst:basecontainer} now provide access to these type IDs.
The first one from \cref{lst:basecontainer1} is a purely virtual method that gets implemented within each derived container class.
The second one (\cref{lst:basecontainer2}) comes in the form of a static method, allowing to check container types using, e.g., \texttt{WIR\_Domination::getContainerTypeID()}.

Finally, the uniqueness of \gls{WIR} containers can be checked using method \texttt{is\-Unique} (\cref{lst:basecontainer3}).
Uniqueness denotes whether at most one instance of a container type can be attached to some \gls{WIR} object.
All containers from \cref{tab:anacont} are unique so that it is not possible to attach, e.g., several domination containers to the very same basic block.
If a new unique container is attached to a \gls{WIR} object using \texttt{insertContainer} (cf. \cref{lst:contapiins}), any possibly already attached containers of the same type are silently erased before.
Class \texttt{WIR\_Comment} (cf. \cref{sec:streammanip}) is a non-unique container so that many comments can be attached to the same \gls{WIR} object.

\begin{lstlisting}[style=code,float=h,caption={Type Handling of Class \texttt{WIR\_BaseContainer}},label=lst:basecontainer]
// Get the unique ID of a container type.
virtual WIR_id_t WIR_BaseContainer::getContainerType@( void )@ const @=@ §0§;Ü\label{lst:basecontainer1}Ü

// Get the unique ID of a container type.
static WIR_id_t WIR_BaseContainer::getContainerType@( void )@;Ü\label{lst:basecontainer2}Ü

// Determine whether a container is unique.
virtual @bool@ WIR_BaseContainer::isUnique@( void )@ const @=@ §0§;Ü\label{lst:basecontainer3}Ü
\end{lstlisting}


%-------------------------------------------------------------------------
\section{Basic Intraprocedural Control Flow Analysis}
\label{sec:cflow}
%-------------------------------------------------------------------------

\begin{tabular}{|l|l|l|l|l|}
  \hline
  \textbf{Analysis Class} & \textbf{Analysis Level} & \textbf{Container Class} & \textbf{Fixing Points} & \textbf{Architectures} \\
  \hhline{|=|=|=|=|=|}
  \texttt{\small WIR\_ControlFlowAnalysis} & Functions & None & None & Generic, all \\
  \hline
\end{tabular}

The foundations for control flow analysis are realized within \gls{WIR} by class \texttt{WIR\_ControlFlowAnalysis} (cf. file \texttt{WIR/analyses/controlflow/wircfganalysis.h}) that actually models a \gls{CFG} according to \cref{def:cfg} on \cpageref{def:cfg}.
This \gls{WIR} analysis automatically generates an internal graph data structure based on a function's basic blocks and their predecessor and successor sets (cf. \cref{lst:cfg}).
This graph is then traversed and the main results of this traversal are various orders in which basic blocks were visited, and classifications of \gls{CFG} nodes and edges~\cite[chapter~7.2]{Much97}.
Since basic blocks and their successor/predecessor relations are modeled in a generic, processor-independent fashion within \gls{WIR}, class \texttt{WIR\_ControlFlowAnalysis} is also generic and does not require any adaptations for specific processor architectures.

\begin{lstlisting}[style=code,float=!b,caption={Public Interface of Class \texttt{WIR\_ControlFlowAnalysis}},label=lst:cflow]
// Default constructor for function-level analysis.
explicit WIR_ControlFlowAnalysis::WIR_ControlFlowAnalysis@(@ WIR_Function æ&æf @)@;Ü\label{lst:cflow1}Ü

// Get the list of visited nodes in depth-first traversal order.
const ästd::listäæ<æästd::reference_wrapperäæ<æWIR_BasicBlockæ>> &æWIR_ControlFlowAnalysis::getDFSOrder@( void )@ const;Ü\label{lst:cflow2}Ü

// Get the list of visited nodes in reverse postorder.
const ästd::listäæ<æästd::reference_wrapperäæ<æWIR_BasicBlockæ>>
  &æWIR_ControlFlowAnalysis::getReversePostOrder@( void )@ const;Ü\label{lst:cflow3}Ü

// Get the list of visited nodes in reverse-topological order.
const ästd::listäæ<æästd::reference_wrapperäæ<æWIR_BasicBlockæ>>
  &æWIR_ControlFlowAnalysis::getReverseTopologicalOrder@( void )@ const;Ü\label{lst:cflow4}Ü

// Get the set of all CFG start nodes.
const WIR_BasicBlockSet æ&æWIR_ControlFlowAnalysis::getStartNodes@( void )@ const;Ü\label{lst:cflow5}Ü

// Get the set of all back edges in the CFG.
const WIR_BasicBlockPairSet æ&æWIR_ControlFlowAnalysis::getBackEdges@( void )@ const;Ü\label{lst:cflow6}Ü

// Visualize a CFG using xdot.
void WIR_ControlFlowAnalysis::visualize@( bool@ i æ=æ false @)@ const;Ü\label{lst:cflow7}Ü
\end{lstlisting}

The public \gls{API} of class \texttt{WIR\_ControlFlowAnalysis} is shown in \cref{lst:cflow}.
This class analyzes the intraprocedural control flow so that caller/callee relationships between functions are not considered.
Due to this, the constructor from \cref{lst:cflow1} expects a reference to a \gls{WIR} function as argument and then builds and traverses the \gls{CFG} for this one function.
\gls{CFG} traversals always start at the very first basic block of a function, as this is a function's unique entry point.
In the case that a function contains unreachable basic blocks, these blocks are also traversed and visited afterwards.

The information produced during these \gls{CFG} traversals is not considered to be ``real'' analysis data providing a higher level of understanding of a \gls{WIR} function.
Instead, it basically serves as basis for other, more sophisticated analyses presented in the following sections.
For this reason, class \texttt{WIR\_ControlFlowAnalysis} does not produce any kind of analysis containers as introduced in \cref{sec:containers}.
Instead, a few getter methods can be used to query \gls{CFG} traversal information.

Methods \texttt{getDFSOrder}, \texttt{getReversePostOrder} and \texttt{getReverseTopologicalOrder} (\crefrange{lst:cflow2}{lst:cflow4}) can then be used to obtain the lists of \gls{CFG} nodes according to the graph's DFS order, postorder or topological order, resp., where the last two orders are reversed in order to support backward analyses.
The set of start nodes of the \gls{CFG}, i.e., basic blocks having no predecessors, is provided by method \texttt{getStartNodes} (\cref{lst:cflow5}).

Back edges inside a \gls{CFG} play an important role, since they close cycles in the graph and thus indicate the presence of loops.
The set of all identified back edges is returned by method \texttt{getBackEdges} (\cref{lst:cflow6}).
\gls{CFG} edges are represented as pairs of basic blocks such that the first element of a pair denotes the edge's source node, and the second element refers to the target node.
The declaration of such basic block pairs can be seen in \cref{lst:bbpairs1} of \cref{lst:bbpairs}.
Based on the concepts introduced in \cref{chap:setsmaps} and by exploiting a dedicated comparator class for such pairs, \cref{lst:bbpairs2} shows the declaration of deterministic \gls{CFG} edge sets.

\begin{lstlisting}[style=code,float=t,caption={Declaration of Basic Block Pairs in file \texttt{wirtypes.h}},label=lst:bbpairs]
// WIR_BasicBlockPair represents pairs of references to WIR basic blocks.
using WIR_BasicBlockPair æ=æ
  ästd::pairäæ<æästd::reference_wrapperäæ<æWIR_BasicBlockæ>æ, ästd::reference_wrapperäæ<æWIR_BasicBlockæ>>æ;Ü\label{lst:bbpairs1}Ü

// WIR_BasicBlockPairSet represents deterministic sets of basic block pairs.
using WIR_BasicBlockPairSet æ=æ ästd::setäæ<æWIR_BasicBlockPair, WIR_Compare_Pairæ>æ;Ü\label{lst:bbpairs2}Ü
\end{lstlisting}

Finally, method \texttt{visualize} from \cref{lst:cflow7} of \cref{lst:cflow} can be used to visualize the created \gls{CFG} for human inspection.
This method converts a \gls{CFG} to the Graphviz DOT language~\cite{dot26}, exports it into a temporary file, and invokes the interactive graph visualization tool \texttt{xdot} on it.
The Boolean argument \texttt{in} of this method defaults to \texttt{false} which means that the internals of \gls{CFG} nodes are not visualized.
By setting this argument explicitly to \texttt{true}, all instructions contained by basic blocks are displayed.

\begin{figure}[h]
  \begin{center}
    \includegraphics[width=0.8\textwidth]{cfg.png}
  \end{center}
  \caption{Control Flow Graph Visualized using Class \texttt{WIR\_ControlFlowAnalysis}}
  \label{fig:cfg}
\end{figure}

\cref{fig:cfg} shows a \gls{CFG} visualization created using the \gls{WIR} unit test from file \texttt{WIR/tests/wirDFS2.cc}.
As can be seen, basic blocks are represented by boxes, where start blocks (here: \texttt{foo} and \texttt{.L5657}) are formatted in bold.
Regular graph edges representing an implicit successorship between two nodes (cf. \cref{sec:basicblocks} on \cpageref{sec:basicblocks}) are displayed in normal style, e.g., the edge from \texttt{.L5654} to \texttt{.L5655}.
So-called true edges are edges where an explicit branch needs to be taken, and they are represented by bold edges (e.g., the edge from \texttt{foo} to \texttt{.L5660} in \cref{fig:cfg}).
Back edges are represented in dashed style, see, e.g., the edge from \texttt{.L5656} to \texttt{.L5654}.


%-------------------------------------------------------------------------
\section{Reachability Analysis}
\label{sec:reachability}
%-------------------------------------------------------------------------

\begin{tabular}{|l|l|l|l|l|}
  \hline
  \textbf{Analysis Class} & \textbf{Analysis Level} & \textbf{Container Class} & \textbf{Fixing Points} & \textbf{Architectures} \\
  \hhline{|=|=|=|=|=|}
  \texttt{\small WIR\_ReachabilityAnalysis} & Functions & \texttt{\small WIR\_Reachability} & Basic Blocks & Generic, all \\
  \hline
\end{tabular}

Reachability analysis as realized by \gls{WIR} class \texttt{WIR\_ReachabilityAnalysis} (cf. file \texttt{WIR/analyses/reach\-ability/wirreachabilityanalysis.h}) aims to answer the question which basic blocks may be reached from a given start block in a \gls{CFG}.
In other words, it determines for each basic block $\texttt{b}_s$ the set of basic blocks $\texttt{b}_t$ such that a directed path from $\texttt{b}_s$ to $\texttt{b}_t$ exists in the \gls{CFG}.
Reachability analysis bases on the generic control flow analysis from \cref{sec:cflow} so that only the intraprocedural \gls{CFG} is considered.
The analysis is entirely generic and does not require adaptations for specific processor architectures.

Reachability analysis is carried out at the level of single \gls{WIR} functions so that the public \gls{API} of class \texttt{WIR\_Reach\-abilityAnalysis} boils down to a function-level constructor as can be seen in \cref{lst:wiranalysis3} of \cref{lst:wiranalysis} and the \texttt{analyze} method from \cref{lst:wiranalysis7}.

Reachability information is attached to basic blocks using container class \texttt{WIR\_Reachability}, cf. file \texttt{WIR/con\-tainers/wirreachability.h}.
A basic block to which a reachability container is attached is supposed to be a start block $\texttt{b}_s$, and the container provides information about all basic blocks $\texttt{b}_t$ reachable from $\texttt{b}_s$.
The relevant parts of its public interface are shown in \cref{lst:reachblocks}.

\begin{lstlisting}[style=code,float=h,caption={Reachability Handling of Container Class \texttt{WIR\_Reachability}},label=lst:reachblocks]
// Add a basic block to the set of reachable blocks.
@void@ WIR_Reachability::addReachableBlock@(@ WIR_BasicBlock æ&æb, @bool@ be æ=æ true @)@;Ü\label{lst:reachblocks1}Ü

// Get the set of all reachable basic blocks.
const WIR_BasicBlockSet æ&æWIR_Reachability::getReachableBlocks@( bool@ be æ=æ true @)@ const;Ü\label{lst:reachblocks2}Ü

// Determine whether a given basic block is reachable from that basic block holding this container.
@bool@ WIR_Reachability::isReachable@(@ const WIR_BasicBlock æ&æb, @bool@ be æ=æ true @)@ const;Ü\label{lst:reachblocks3}Ü

// Determine whether the latter instruction is reachable from the former.
static @bool@ WIR_Reachability::isReachable@(@ const WIR_Instruction æ&æi1, const WIR_Instruction æ&æi2,                                                     @bool@ be æ=æ true @)@ const;Ü\label{lst:reachblocks4}Ü
\end{lstlisting}

A basic block can be added to the set of reachable blocks using method \texttt{addReachableBlock} (cf. \cref{lst:reachblocks1}).
This method (as well as all others shown in \cref{lst:reachblocks}) accepts a Boolean flag defaulting to \texttt{true} as optional last argument.
This flag indicates whether reachability across back edges in the \gls{CFG} shall be considered or not, i.e., whether the analyzed paths from $\texttt{b}_s$ to $\texttt{b}_t$ are allowed to include back edges or not.
The set of all blocks $\texttt{b}_t$ that are reachable from the block hosting the container can be obtained via method \texttt{getReachableBlocks} (\cref{lst:reachblocks2}).
The variant of the \texttt{isReachable} method from \cref{lst:reachblocks3} can be used to query whether a given basic block is part of the set of reachable blocks.
A static variant of \texttt{isReachable} exists (cf. \cref{lst:reachblocks4}) in order to simply check reachability between two given instructions.
This method takes the basic blocks holding the given instructions and checks the reachability of these blocks.
If both instructions reside in the same basic block, \texttt{isReachable} considers which of the two instructions comes first in that basic block.

\clearpage


%-------------------------------------------------------------------------
\section{Domination and Immediate Domination Analysis}
\label{sec:domination}
%-------------------------------------------------------------------------

\begin{tabular}{|l|l|l|l|l|}
  \hline
  \textbf{Analysis Class} & \textbf{Analysis Level} & \textbf{Container Class} & \textbf{Fixing Points} & \textbf{Architectures} \\
  \hhline{|=|=|=|=|=|}
  \texttt{\small WIR\_DominationAnalysis} & Functions & \texttt{\small WIR\_Domination} & Basic Blocks & \multirow{4}{*}{Generic, all} \\
  & & & Instructions & \\
  \hhline{----}
  \texttt{\small WIR\_ImmediateDomination-} & Functions & \texttt{\small WIR\_Immediate-} & Basic Blocks & \\
  \texttt{\small \ \ Analysis} & & \texttt{\small \ \ Domination} & & \\
  \hline
\end{tabular}

The domination and immediate domination analyses within \gls{WIR} are realized by classes \texttt{WIR\_DominationAnaly\-sis} (cf. file \texttt{WIR/analyses/domination/wirdominationanalysis.h}) and \texttt{WIR\_ImmediateDomination\-Analysis} (file \texttt{WIR/analyses/immediatedomination/wirimmediatedominationanalysis.h}).
According to~\cite[chapter~7.3]{Much97}, domination is a relation \emph{dom} between two \gls{CFG} nodes \texttt{d} and \texttt{n} such that
\begin{equation}
  \texttt{d}\ \textit{dom}\ \texttt{n} \Leftrightarrow \text{Every path from the CFG's start node to \texttt{n} includes \texttt{d}}
\end{equation}
Immediate domination is a sub-relation \emph{idom} such that
\begin{equation}
  \forall\texttt{i} \not= \texttt{n}:\ \texttt{i}\ \textit{idom}\ \texttt{n} \Leftrightarrow \texttt{i}\ \textit{dom}\ \texttt{n}\ \wedge \nexists\texttt{c}:\ (\texttt{c} \not= \texttt{i}) \wedge (\texttt{c} \not= \texttt{n}) \wedge (\texttt{i}\ \textit{dom}\ \texttt{c}) \wedge (\texttt{c}\ \textit{dom}\ \texttt{n})
\end{equation}
Nodes \texttt{d} that (immediately) dominate some \texttt{n} are called (immediate) dominators.
As was the case in the previous sections, (immediate) domination analysis takes place at the function-level, and both analyses are entirely generic.
Their implementations closely follow the algorithms from~\cite[chapter~18.1]{Appel04} and~\cite[chapter~7.3]{Much97}.

Usually, domination relates to individual basic blocks as \gls{CFG} nodes \texttt{d}, \texttt{i} or \texttt{n} in the above equations.
Thus, analysis results are regularly attached to basic blocks using container classes \texttt{WIR\_Domination} and \texttt{WIR\_Immediate\-Domination}, resp.
For domination analysis, however, it might be useful to break the basic block-level relation \textit{dom} down also to the level of single \gls{WIR} instructions. Then, \texttt{WIR\_Domination} containers are attached to both basic blocks and instructions.
This behavior of class \texttt{WIR\_DominationAnalysis} can be controlled using its constructor shown in \cref{lst:doma1} of \cref{lst:doma}.
Besides the function to be analyzed, it obtains a Boolean defaulting to \texttt{false} that specifies whether instruction-level domination shall be computed.
Except this Boolean flag, the \gls{API} of class \texttt{WIR\_ImmediateDominationAnalysis} is identical to \cref{lst:doma}.

\begin{lstlisting}[style=code,float=h,caption={Public Interface of Class \texttt{WIR\_DominationAnalysis}},label=lst:doma]
// Default constructor for function-level analysis.
explicit WIR_DominationAnalysis::WIR_DominationAnalysis@(@ WIR_Function æ&æf, @bool@ i æ=æ false @)@;Ü\label{lst:doma1}Ü

// Performs the actual code analysis.
@void@ WIR_DominationAnalysis::analyze@( void )@;Ü\label{lst:doma2}Ü
\end{lstlisting}

\cref{lst:dom} shows the interfaces for (immediate) domination handling of both container classes.
Dominators in the form of basic blocks or instructions can be added to class \texttt{WIR\_Domination} using the \texttt{addDominator} methods from \cref{lst:dom1} or \cref{lst:dom2}, resp.
Likewise, either the set of dominator blocks or instructions can be retrieved using \texttt{getDominatorBlocks} (\cref{lst:dom3}) or \texttt{getDominatorInstructions} (\cref{lst:dom4}).
Relation \textit{idom} is unique so that a basic block \texttt{n} has exactly one immediate dominator \texttt{i}.
For this reason, container class \texttt{WIR\_ImmediateDomination} simply handles a reference to one single immediate dominator block.
Method \texttt{setImmediateDominator} (\cref{lst:dom5}) can be used to set this immediate dominator, while \texttt{getImmediateDominator} (\cref{lst:dom6}) retrieves the immediate dominator.

\begin{lstlisting}[style=code,float=ht,caption={Domination Handling of Container Classes \texttt{WIR\_Domination} and \texttt{WIR\_ImmediateDomination}},label=lst:dom]
// Add a basic block to the set of dominators.
@void@ WIR_Domination::addDominator@(@ const WIR_BasicBlock æ&æb @)@;Ü\label{lst:dom1}Ü

// Add an instruction to the set of dominators.
@void@ WIR_Domination::addDominator@(@ const WIR_Instruction æ&æi @)@;Ü\label{lst:dom2}Ü

// Get the set of all dominator basic blocks.
const WIR_BasicBlockSet æ&æWIR_Domination::getDominatorBlocks@( void )@ const;Ü\label{lst:dom3}Ü

// Get the set of all dominator instructions.
const WIR_InstructionSet æ&æWIR_Domination::getDominatorInstructions@( void )@ const;Ü\label{lst:dom4}Ü


// Set a basic block as immediate dominator.
@void@ WIR_ImmediateDomination::setImmediateDominator@(@ const WIR_BasicBlock æ&æb @)@;Ü\label{lst:dom5}Ü

// Get the immediate dominator basic block.
const WIR_BasicBlock æ&æWIR_Domination::getImmediateDominator@( void )@ const;Ü\label{lst:dom6}Ü
\end{lstlisting}


%-------------------------------------------------------------------------
\section{Structural Control Flow Analysis}
\label{sec:structcflow}
%-------------------------------------------------------------------------

\begin{tabular}{|l|l|l|l|l|}
  \hline
  \textbf{Analysis Class} & \textbf{Analysis Level} & \textbf{Container Class} & \textbf{Fixing Points} & \textbf{Architectures} \\
  \hhline{|=|=|=|=|=|}
  \texttt{\small WIR\_StructuralAnalysis} & Functions & \texttt{\small WIR\_ControlTree} & Basic Blocks & Generic, all \\
  \hline
\end{tabular}

The goal of structural control flow analysis is to decompose a general \gls{CFG} of a function into (possibly nested) hierarchical regions that provide more high-level, structural control flow information.
Regularly, the information provided by a \gls{CFG} only models predecessor and successor relations between nodes.
It remains unclear at this level, which groups of nodes form, e.g., an if-else construct or a tail-controlled loop.
This changes during structural control flow analysis, where typical control flow structures well-known from high-level programming languages are systematically identified.

This analysis is realized by \gls{WIR} class \texttt{WIR\_StructuralAnalysis} declared in file \texttt{WIR/analyses/structural\-controlflow/wirstructuralanalysis.h}.
It closely follows the concepts presented in~\cite[chapter~7.7]{Much97}.
In analogy to all control flow analyses presented before, structural analysis is performed at the level of single functions, too.
Furthermore, it requires no particular processor-specific extensions.

\begin{lstlisting}[style=code,float=!b,caption={Public Interface of Class \texttt{WIR\_StructuralAnalysis}},label=lst:structcflow]
// Default constructor for function-level analysis.
explicit WIR_StructuralAnalysis::WIR_StructuralAnalysis@(@ WIR_Function æ&æf, @bool@ verbose æ=æ false,                                                                    @bool@ tmpFiles æ=æ false @)@;Ü\label{lst:structcflow1}Ü

// Performs the actual code analysis.
@void@ WIR_StructuralAnalysis::analyze@( void )@;Ü\label{lst:structcflow2}Ü

// Visualize the hierarchical CFG using xdot.
void WIR_StructuralAnalysis::visualize@( void )@ const;Ü\label{lst:structcflow3}Ü
\end{lstlisting}

\cref{lst:structcflow} shows the \gls{API} of class \texttt{WIR\_StructuralAnalysis}.
Besides the \gls{WIR} function to be analyzed, its constructor (cf. \cref{lst:structcflow1}) features two optional Boolean arguments that steer some technical aspects of method \texttt{visualize} from \cref{lst:structcflow3}.
Similar to \cref{sec:cflow}, this method serves to visualize the outcomes of structural analysis for a human inspection.
This again is done by plotting the hierarchical \gls{CFG} in the external tool \texttt{xdot}.
The constructor's parameter \texttt{verbose} controls whether the command lines used to invoke \texttt{xdot} are printed.
Flag \texttt{tmpFiles} has the effect that the temporary DOT file that gets displayed using \texttt{xdot} is not erased.
All subsequent figures illustrating the various kinds of control regions identified during structural analysis have been generated using the \texttt{visualize} method.

All these regions have in common that they have exactly one entry point.
This property has the effect that any arbitrary unstructured \gls{CFG} becomes reducible.
Reducibility means that a series of transformations can be applied to the initial \gls{CFG} clustering individual nodes into hierarchical nodes so that, in the end, the entire \gls{CFG} gets reduced to a single (hierarchical) node~\cite[chapter~7.5]{Much97}.
Due to this property, the result produced by structural analysis no longer is a general graph, but a tree instead.
The root node of this hierarchical control tree represents the entire function that has been analyzed.
The tree's leaf nodes represent all the individual basic blocks of the original, unstructured \gls{CFG}.
In the following, we will use the terms ``control region'' and ``tree node'' interchangably for this reason.

\begin{figure}[t]
  \begin{center}
    \includegraphics[width=0.63\textwidth]{controltreenodes.png}
  \end{center}
  \caption{Inheritance Diagram of Control Regions}
  \label{fig:ctnodes}
\end{figure}

All the different control regions identified during structural analysis are represented by \gls{WIR} classes that inherit from the abstract base class \texttt{WIR\_ControlTreeNode} as shown in \cref{fig:ctnodes}.
The general properties of this base class are shown in \cref{lst:ctnode}.
As usual within \gls{WIR}, there is a \texttt{getType} method (cf. \cref{lst:ctnode1}) that allows to conveniently determine a region's type using the enumeration type \texttt{WIR\_CTNodeType} from \cref{lst:ctnode11}.
Structural analysis generally distinguishes between acyclic and cyclic control regions, and these properties can be queried using methods \texttt{isAcyclic} (\cref{lst:ctnode2}) and \texttt{isCyclic} (\cref{lst:ctnode3}), resp.
A region's abovementioned unique entry can be obtained with method \texttt{getEntry} (\cref{lst:ctnode4}) that returns a constant reference to the identified entry node.
Whether or not a region actually is some entry of another, surrounding region can be checked using \texttt{isEntry} (\cref{lst:ctnode5}).

\begin{lstlisting}[style=code,float=!t,caption={Public Interface of Class \texttt{WIR\_ControlTreeNode}},label=lst:ctnode]
// Get a control region's type.
virtual WIR_CTNodeType WIR_ControlTreeNode::getType@( void )@ const æ=æ §0§;Ü\label{lst:ctnode1}Ü

// Determine whther a control region is acyclic or not.
@bool@ WIR_ControlTreeNode::isAcyclic@( void )@ const;Ü\label{lst:ctnode2}Ü

// Determine whther a control region is cyclic or not.
virtual @bool@ WIR_ControlTreeNode::isCyclic@( void )@ const æ=æ §0§;Ü\label{lst:ctnode3}Ü

// Get a control region's unique entry node.
virtual const WIR_ControlTreeNode æ&æWIR_ControlTreeNode::getEntry@( void )@ const æ=æ §0§;Ü\label{lst:ctnode4}Ü

// Determine whether a region is the entry node of some parent or not.
@bool@ WIR_ControlTreeNode::isEntry@( void )@ const;Ü\label{lst:ctnode5}Ü

// Get the set of all child regions.
const WIR_ControlTreeNodeSet æ&æWIR_ControlTreeNode::getChilds@( void )@ const;Ü\label{lst:ctnode6}Ü

// Get a reference to a region's immediate parent in the hierarchy, or a reference to itself.
WIR_ControlTreeNode æ&æWIR_ControlTreeNode::getParent@( void )@ const;Ü\label{lst:ctnode7}Ü

// Get a reference to a region's top-level root in the hierarchy, or a reference to itself.
WIR_ControlTreeNode æ&æWIR_ControlTreeNode::getRoot@( void )@ const;Ü\label{lst:ctnode8}Ü

// Get the set of all basic blocks in the hierarchy below this regions.
WIR_BasicBlockSet WIR_ControlTreeNode::getBasicBlocks@( void )@ const;Ü\label{lst:ctnode9}Ü

// Get a region's set of back edges.
const WIR_ControlTreeEdgeSet æ&æWIR_ControlTreeNode::getBackEdges@( void )@ const;Ü\label{lst:ctnode10}Ü

// Enum WIR_CTNodeType represents different types of WIR control regions.
enum class WIR_CTNodeType : @char@Ü\label{lst:ctnode11}Ü
{
  bb,           // A single basic block.
  block,        // A block schema.
  ifthen,       // An if-then region.
  ifthenelse,   // An if-then-else region.
  switchcase,   // A switch-case region.
  selfloop,     // A self loop.
  whileloop,    // A regular for- or while-do loop.
  naturalloop,  // A natural do-while loop.
  proper,       // An acyclic proper region.
  improper      // A cyclic improper region.
};
\end{lstlisting}

The next methods \texttt{getChilds} (\cref{lst:ctnode6}), \texttt{getParent} (\cref{lst:ctnode7}) and \texttt{getRoot} (\cref{lst:ctnode8}) permit to navigate within the hierarchy of nested control regions in analogy to hierarchical registers (cf. \cref{sssec:regs}).
Method \texttt{getBasic\-Blocks} (\cref{lst:ctnode9}) provides the set of (wrapped) references to all basic blocks that are somehow (hierarchically) contained by a control region.
For cyclic regions, \texttt{getBackEdges} (\cref{lst:ctnode10}) returns a reference to the set of identified back edges, and the returned type \texttt{WIR\_ControlTreeEdgeSet} is declared similar to the standard \gls{CFG} edges from \cref{lst:bbpairs}.

\cref{fig:acyclic} shows the acyclic control structures that are recognized by the analysis.
We quickly introduce the meaning and the key properties of these regions in the following.

\begin{description}
  \item[Basic blocks:] Basic block tree nodes form the leafs of the hierarchical control tree produced by structural analysis.
  In \cref{fig:acyclic}, they are represented by the ellipses labeled with the block names.
  Due to their simplicity, class \texttt{WIR\_BasicBlockTreeNode} provides no particular interface except method \texttt{getBasicBlock} that allows to retrieve the actual \gls{WIR} basic block:
\begin{lstlisting}[style=code,float=!h,xleftmargin=10mm]
// Get the basic block represented by this region.
WIR_BasicBlock æ&æWIR_BasicBlockTreeNode::getBasicBlock@( void )@ const;
\end{lstlisting}

  \item[Block regions:] Block regions are not to be confused with the basic block nodes above, as block regions represent sequences of maximal length of tree nodes.
  Control flow enters a block region at its entry node and only follows the sequence of nodes contained by the block region.
  When reaching the block region's very last node, control flow leaves the region.
  Due to this sequential behavior, block regions have both the single-entry and single-exit property, see \cref{fig:acyclic}.
  Access to the sequence of nodes contained by a block region is provided using method \texttt{getBlockList} that returns the list of (wrapped) references to the region's nodes in their sequential order:
\begin{lstlisting}[style=code,float=!h,xleftmargin=10mm]
// Get the list of all contained nodes in their sequential order.
const ästd::listäæ<æästd::reference_wrapperäæ<æWIR_ControlTreeNodeæ>> &æWIR_BlockTreeNode::getBlockList@( void )@ const;
\end{lstlisting}

\begin{figure}[t]
  \begin{center}
    \includegraphics[width=0.86\textwidth]{acyclicregions.png}
  \end{center}
  \caption{Acyclic, Structured Control Regions}
  \label{fig:acyclic}
\end{figure}

  \item[If-then structures:] An if-then region contains exactly two nodes, a condition node and a branch node.
  The condition node is the region's entry and depending on its evaluation of some condition, control flow proceeds either to the branch node or to some other node outside the if-then.
  All involved nodes (condition, branch and outside node) are different, the condition node has exactly two successors, and the branch node has exactly one predecessor (the branch node) and one successor.
  Both the condition and the branch node share the same node outside the region as successor (cf. \cref{fig:acyclic}).
  The region's condition and branch nodes can be accessed using the below methods of class \texttt{WIR\_IfThenTreeNode}:
\begin{lstlisting}[style=code,float=!h,xleftmargin=10mm]
// Get an if-then's condition node.
const WIR_ControlTreeNode æ&æWIR_IfThenTreeNode::getCondition@( void )@ const;

// Get an if-then's then-branch node.
const WIR_ControlTreeNode æ&æWIR_IfThenTreeNode::getBranch@( void )@ const;
\end{lstlisting}

  \item[If-then-else structures:] Similar to if-then regions, the if-then-else region consists of a condition node branching to two different successors.
  However, these successors represent the two different branches of an if-then-else construct and thus are both part of the region.
  The two branches continue with the very same exit node outside the if-then-else, see \cref{fig:acyclic}.
  As a consequence, class \texttt{WIR\_IfThenElseTreeNode} provides the methods \texttt{getCondition}, \texttt{getBranch1} and \texttt{Branch2} to access its internal nodes:
\begin{lstlisting}[style=code,float=!ht,xleftmargin=10mm]
// Get an if-then-else's condition node.
const WIR_ControlTreeNode æ&æWIR_IfThenElseTreeNode::getCondition@( void )@ const;

// Get an if-then-else's first branch node.
const WIR_ControlTreeNode æ&æWIR_IfThenElseTreeNode::getBranch1@( void )@ const;

// Get an if-then-else's second branch node.
const WIR_ControlTreeNode æ&æWIR_IfThenElseTreeNode::getBranch2@( void )@ const;
\end{lstlisting}

  \item[Switch-cases:] Are a generalization of if-then-else structures in that sense that the condition node can branch to more than 2 cases within this region, cf. \cref{fig:acyclic}.
  All the case nodes share the condition node as predecessor.
  A case node can break out of the switch-case to a node outside the region as successor, where all cases share this very same node as common exit node.
  Alternatively, a case node can proceed to one other case node in the presence of a fall-through (not shown in \cref{fig:acyclic}).
  All predecessors of case nodes must be either the condition node or some other case node.
  The public \gls{API} of class \texttt{WIR\_SwitchCaseTreeNode} is quite straightforward, methods \texttt{getCondition} and \texttt{getCases} below return the region's condition node and the list of all identified case nodes, resp.
  Method \texttt{getFallthroughCases} helps to identify all those case nodes that exit with a fallthrough.
  It returns a list of iterators, where each iterator refers to the position of a fallthrough case in the list provided by \texttt{getCases}.
\begin{lstlisting}[style=code,float=!h,xleftmargin=10mm]
// Get a switch-case's condition node.
const WIR_ControlTreeNode æ&æWIR_SwitchCaseTreeNode::getCondition@( void )@ const;

// Get the list of all contained case nodes in their sequential order.
const ästd::listäæ<æästd::reference_wrapperäæ<æWIR_ControlTreeNodeæ>>                                        &æWIR_SwitchCaseTreeNode::getCases@( void )@ const;

// Get the list of all fall-through case nodes in their sequential order.
ästd::listäæ<æästd::listäæ<æästd::reference_wrapperäæ<æWIR_ControlTreeNodeæ>>::const_iterator>æ WIR_SwitchCaseTreeNode::getFallthroughCases@( void )@ const;
\end{lstlisting}
\end{description}

All control regions presented so far and shown in \cref{fig:acyclic} are acyclic regions that do not possess any back edges.
\cref{fig:cyclic} now depicts the regular cyclic regions recognized by the structural analysis.

\begin{figure}[h]
  \begin{center}
    \includegraphics[width=0.6\textwidth]{cyclicregions.png}
  \end{center}
  \caption{Cyclic, Structured Control Regions}
  \label{fig:cyclic}
\end{figure}

\begin{description}
  \item[Self-loops:] The most simple cyclic control region is a self-loop that consists of exactly one node that either branches back to itself or to some outside node (cf. \cref{fig:cyclic}).
  The corresponding class \texttt{WIR\_SelfLoop\-TreeNode} features no particular methods that go beyond that of its base class shown in \cref{lst:ctnode}.
  Especially, method \texttt{getEntry} (cf. \cref{lst:ctnode4} of \cref{lst:ctnode}) can be used to retrieve a self-loop's only node.

  \item[While-do loops:] This region represents classical head-controlled loops like, e.g., for or while-do loops.
  A head-controlled loop is characterized by an entry node that checks the loop's condition and either continues with the loop body or some node outside the loop region.
  Thus, this control region consists of two nodes, the entry node and the loop body.
  Only predecessor of the loop body is the entry node, and the loop body jumps back to the entry node via a back edge.
  The body can break out of its loop via break edges (not shown in \cref{fig:acyclic}).
  In this case, the target of such a break edge must be the same node outside the loop that is also reached from the loop's entry.
  Again, the region's class \texttt{WIR\_WhileLoopTreeNode} is simple -- besides method \texttt{getEntry} to obtain the loop's entry node, it only comes with method \texttt{getNodes} returning the set of all nodes in this region:
\begin{lstlisting}[style=code,float=!h,xleftmargin=10mm]
// Get a while loop's inner nodes.
const WIR_ControlTreeNodeSet æ&æWIR_WhileLoopTreeNode::getNodes@( void )@ const;
\end{lstlisting}

  The returned set of type \texttt{WIR\_ControlTreeNodeSet} is defined in analogy to the concepts discussed in \cref{chap:setsmaps}.

  \item[Natural loops:] Represent tail-controlled loops like, e.g., do-while loops, where the loop's exit condition is checked after the execution of the loop body.
  This implies that the region's entry node represents the loop body, and the body's successor is the loop's condition node.
  This condition node either branches back to the body, or to some exit node outside the loop (cf. \cref{fig:cyclic}).
  Class \texttt{WIR\_NaturalLoopTreeNode} has exactly the same public \gls{API} as while-do loops have so that we omit its discussion for the sake of brevity.
\end{description}

The loop regions shown in \cref{fig:cyclic} all have a unique single entry point, but can have multiple exits in the case of breaks.
For a structural control flow analysis, it is difficult to unambiguously distinguish the loop's regular exit edge from potential break edges.
In order to support the analysis, so-called break containers of class \texttt{WIR\_Break} (cf. file \texttt{WIR/containers/wirbreak.h}) can be attached to branch operations that jump to a target outside a loop.
If, e.g., a compiler's code selector creates this kind of hints whenever it exhibits a loop break, \gls{WIR}'s structural analysis can produce loop regions that precisely match the source code's structure.
The container class \texttt{WIR\_Break} comes with absolutely no specific properties or methods, the structural analysis simply checks whether such a container exists for a given branch or not.

All regions shown in \cref{fig:acyclic,fig:cyclic} are well-structured in that sense that they inherently possess the single-entry property and thus lead to \gls{CFG} reducibility.
However, certain programming languages like, e.g., ANSI-C, permit to generate unstructured code where certain regions can be entered via several nodes.
Clearly, such unstructured constructs should be avoided whenever possible, but structural analysis needs to be able to cope with such unstructured regions.
For this reason, the analysis features two additional control regions which are depicted in \cref{fig:unstructured}.

\begin{figure}[t]
  \begin{center}
    \includegraphics[width=0.6\textwidth]{unstructured.png}
  \end{center}
  \caption{Unstructured Control Regions}
  \label{fig:unstructured}
\end{figure}

\begin{description}
  \item[Proper regions:] Such regions are defined to be acyclic, single-entry regions of minimal size that cannot be reduced to any of the well-structured acyclic regions from \cref{fig:acyclic}.
  When considering only the basic blocks \texttt{.L2}, \texttt{.L4}, \texttt{.L5} and \texttt{.L6} from \cref{fig:unstructured}, then this part of the \gls{CFG} resembles an if-then-else structure.
  However, branch \texttt{.L5} can be reached via a second entry node, \texttt{.L3}.
  Thus, there are several entries into this unstructured ``if-then-else'' construct.
  In such a case, \gls{WIR}'s structural analysis still enforces that only single-entry regions are created.
  Thus, the common dominator node (cf. \cref{sec:domination}) that is closest to the multiple entries (here: \texttt{.L2} and \texttt{.L3}) is determined and taken as single entry node for the created unstructured, proper region shown in \cref{fig:unstructured}.
  The nodes contained by a proper region can be retrieved using methods \texttt{getEntry} and \texttt{getNodes} below:
\begin{lstlisting}[style=code,float=!h,xleftmargin=10mm]
// Get a proper region's entry node.
virtual const WIR_ControlTreeNode æ&æWIR_ProperTreeNode::getEntry@( void )@ const;

// Get a proper region's inner nodes.
const WIR_ControlTreeNodeSet æ&æWIR_ProperTreeNode::getNodes@( void )@ const;
\end{lstlisting}

  \item[Improper regions:] Are the same as proper regions, but just for cyclic unstructured  \gls{CFG} regions.
  \cref{fig:unstructured} shows an example of a while-do loop consisting of the entry \texttt{.L8} and the loop body \texttt{.L9}.
  However, the loop body can also be reached from some other node \texttt{.L7} outside the loop so that this code structure cannot be reduced to any of the regular structures from \cref{fig:cyclic}.
  Improper regions are represented by \gls{WIR} class \texttt{WIR\_ImproperTreeNode}, its interface is identical to that of proper regions above.
\end{description}


%-------------------------------------------------------------------------
\section{Flow Facts for Execution Time Analysis}
\label{sec:flowfacts}
%-------------------------------------------------------------------------

\begin{tabular}{|l|l|l|l|l|}
  \hline
  \textbf{Analysis Class} & \textbf{Analysis Level} & \textbf{Container Class} & \textbf{Fixing Points} & \textbf{Architectures} \\
  \hhline{|=|=|=|=|=|}
  Does not apply &  Does not apply & \texttt{\small WIR\_FlowFactRef} & Basic Blocks & Generic, all \\
  \hline
\end{tabular}

Static execution time analysis aims to estimate the actual execution time of a piece of code.
The analysis' results usually are the numbers of clock cycles or the physical runtimes of \gls{CFG} nodes, functions, or entire tasks.
Such an execution time analysis regularly inspects the set of possible paths through a task's \gls{CFG}.
Especially in the presence of cyclic code structures, the number of iterations through each cycle needs to be estimated.
This, however, is impossible in general due to the undecidability of the halting problem.
For this reason, static timing analyses require human-provided meta-information giving hints about how often control regions are executed.
These hints are called flow facts.

\begin{figure}[t]
  \begin{center}
    \includegraphics[width=0.86\textwidth]{flowfacts.png}
  \end{center}
  \caption{Inheritance Diagram of Flow Facts}
  \label{fig:ffacts}
\end{figure}

\gls{WIR} does not feature a static timing analysis by itself.
However, it contains mechanisms for the annotation and query of flow facts, and this section describes all the underlying concepts.
All involved flow fact classes can be found in folder \texttt{WIR/flowfacts}, and they are all generic and processor-independent.

\gls{WIR} currently supports three different kinds of flow facts: entry points, loop bounds, and flow restrictions, which all inherit from the abstract base class \texttt{WIR\_FlowFact} (cf. file \texttt{wirflowfact.h}).
\cref{fig:ffacts} shows the inheritance diagram of all flow fact classes.
The most relevant properties of this abstract base class are summarized in \cref{lst:ffact}.
As usual, a method \texttt{getType} (cf. \cref{lst:ffact1}) can be used to quickly determine the type of a flow fact.
The result is provided by enumeration type \texttt{WIR\_FlowFactType} from \cref{lst:ffact5}.
There also exists an output \texttt{<{}<} operator in order to print out flow facts (\cref{lst:ffact2}).
Flow facts are centrally managed by a \gls{WIR} system.
Thus, method \texttt{isInserted} (\cref{lst:ffact3}) can be used to determine whether a flow fact is properly inserted into a system.
If so, the system associated with a flow fact can be obtained using \texttt{getSystem} (\cref{lst:ffact4}).

\begin{lstlisting}[style=code,float=t,caption={Public Interface of Class \texttt{WIR\_FlowFact}},label=lst:ffact]
// Get a flow fact's type.
virtual WIR_FlowFactType WIR_FlowFact::getType@( void )@ const æ=æ §0§;Ü\label{lst:ffact1}Ü

// Dump a flow fact to an output stream in a human-readable way.
friend ästd::ostreamä æ&æ operator æ<<æ @(@ ästd::ostreamä æ&æos, const WIR_FlowFact æ&æff @)@;Ü\label{lst:ffact2}Ü

// Determine whether a flow fact is inserted into some system.
@bool@ WIR_FlowFact::isInserted@( void )@ const;Ü\label{lst:ffact3}Ü

// Get the system containing a flow fact.
WIR_System æ&æWIR_FlowFact::getSystem@( void )@ const;Ü\label{lst:ffact4}Ü

// Enum WIR_FlowFactType represents different types of WIR flow facts.
enum class WIR_FlowFactType : @char@Ü\label{lst:ffact5}Ü
{
  entrypoint,       // An entry point.
  flowrestriction,  // A flow restriction.
  loopbound         // A loop bound.
};
\end{lstlisting}

\begin{lstlisting}[style=code,float=!b,caption={Flow Fact Handling of Class \texttt{WIR\_System}},label=lst:sysffacts]
// Copy-add a new flow fact at the end of the list of flow facts.
WIR_FlowFact æ&æWIR_System::pushBackFlowFact@(@ const WIR_FlowFact æ&æff @)@;Ü\label{lst:sysffacts1}Ü

// Get the complete list of references of all attached flow facts.
const ästd::listäæ<æästd::reference_wrapperäæ<æWIR_FlowFactæ>> &æWIR_System::getFlowFacts@( void )@ const;Ü\label{lst:sysffacts2}Ü

// Get the set of all flow facts of a specified type.
templateæ<ætypename Tæ>æ const ästd::listäæ<æästd::reference_wrapperäæ<æTæ>>æ WIR_System::getFlowFacts@( void )@ const;Ü\label{lst:sysffacts3}Ü

// Determine whether a system contains a flow fact with the given ID.
@bool@ WIR_System::containsFlowFact@(@ WIR_id_t id @)@ const;Ü\label{lst:sysffacts4}Ü

// Remove the flow fact from the specified list position.
@void@ WIR_System::eraseFlowFact@(@ ästd::listäæ<æästd::reference_wrapperäæ<æWIR_FlowFactæ>>::const_iteratoræ pos @)@;Ü\label{lst:sysffacts5}Ü

// Remove the flow fact with the given ID.
@void@ WIR_System::eraseFlowFact@(@ WIR_id_t id @)@;Ü\label{lst:sysffacts6}Ü

// Remove all flow facts.
@void@ WIR_System::clearFlowFacts@( void )@;Ü\label{lst:sysffacts7}Ü
\end{lstlisting}

Before going into the details of the various classes of flow facts and their semantics, the public interface of class \texttt{WIR\_System} for flow fact management is shown in \cref{lst:sysffacts}.
Flow facts can be added to a system using \texttt{pushBackFlowFact} (cf. \cref{lst:sysffacts1}).
This method returns a reference to the newly added flow fact.
A list with (wrapped) references to all flow facts of a system can be obtained using \texttt{getFlowFacts} from \cref{lst:sysffacts2}.
The list of already correctly casted flow facts of a specific type can be retrieved by the variant of method \texttt{getFlowFacts} from \cref{lst:sysffacts3}.
\texttt{containsFlowFact} (\cref{lst:sysffacts4}) returns whether a flow fact with a specific numerical ID exists within a system or not.
The first variant of method \texttt{eraseFlowFact} from \cref{lst:sysffacts5} erases the flow fact to which a given iterator in the list of all flow facts refers.
The second \texttt{eraseFlowFact} variant (\cref{lst:sysffacts6}) erases the flow fact with a given ID.
Finally, method \texttt{clearFlowFacts} (\cref{lst:sysffacts7}) erases all flow facts attached to a system.
As usual within \gls{WIR}, erased flow facts are automatically destroyed and removed from memory.

The most simple class of flow facts are the entry points represented by class \texttt{WIR\_EntryPoint}.
Their purpose is to mark a given function as entry of a software task modeled within \gls{WIR}.
For single-task systems created from, e.g., C code, the one and only entry point is the function named \texttt{main}.
For multi-task systems, however, there can be various functions that get executed by a system's scheduler.
These functions represent task entries and thus should be identified using entry point flow facts.

\begin{wirNote}
  As already mentioned on \cpageref{futurework:taskmanager}, \gls{WIR} currently does not contain any well-designed system-level task model.
  This will probably change in the future.
  Since entry point flow facts inherently address task-level information, it is very likely that the entire class \texttt{WIR\_EntryPoint} will become obsolete since it gets subsumed into the future task managers mentioned on \cpageref{futurework:taskmanager}.
  Thus, the \gls{API} of entry points discussed in the following should be considered unstable/obsolete and should not actively be used if possible.
\end{wirNote}

\begin{lstlisting}[style=code,float=t,caption={Public Interface of Class \texttt{WIR\_EntryPoint}},label=lst:entrypoint]
// Default constructor.
WIR_EntryPoint::WIR_EntryPoint@(@ const WIR_Function æ&æf, const ästd::mapäæ<æästd::stringä, ästd::stringäæ> &æa );Ü\label{lst:entrypoint1}Ü

// Copy constructor.
WIR_EntryPoint::WIR_EntryPoint@(@ const WIR_EntryPoint æ&æep @)@;Ü\label{lst:entrypoint2}Ü

// Set the function this entry point describes.
@void@ WIR_EntryPoint::setFunction@(@ const WIR_Function æ&æf @)@;Ü\label{lst:entrypoint3}Ü

// Get the function marked as entry point.
WIR_Function æ&æWIR_EntryPoint::getFunction@( void )@ const;Ü\label{lst:entrypoint4}Ü

// Get the task attributes associated with an entry point.
const ästd::mapäæ<æästd::stringä, ästd::stringäæ> &æWIR_EntryPoint::getAttributes@( void )@ const;Ü\label{lst:entrypoint5}Ü
\end{lstlisting}

\cref{lst:entrypoint} shows the relevant parts of the interface of class \texttt{WIR\_EntryPoint}.
As entry points aim to mark \gls{WIR} functions as task entries, the constructor from \cref{lst:entrypoint1} expects a function being task entry as first argument.
The map being the constructor's second argument intends to specify attributes of a task like, e.g., periods or deadlines.
This attribute map is fully unstructured as it only uses plain strings (e.g., key \texttt{"period"} could be mapped to value \texttt{"50ms"}), and such task attributes are currently not used and supported within \gls{WIR} (see above note).
Thus, it is sufficient to just pass the empty map to this constructor.
\cref{lst:entrypoint2} shows that the class comes with a copy constructor, too.
The function marked as entry point can explicitly be set using \texttt{setFunction} (\cref{lst:entrypoint3}) and retrieved via \texttt{getFunction} (\cref{lst:entrypoint4}).
The (useless) attribute map could be obtained with method \texttt{getAttributes} from \cref{lst:entrypoint5}.

Loop bound flow facts intend to specify the minimum and maximum number of times that a loop iterates, as illustrated in the following code examples.
\cref{lst:lb1} shows a loop with a constant number of iterations that could be annotated with a loop bound flow fact carrying exactly this iteration count both as minimum and maximum value.
A data dependent loop whose number of iterations can vary but still is known to be bounded can be represented by a loop bound as can be seen in \cref{lst:lb2}.

\begin{lstlisting}[style=code,float=!h]
for ( i = §100§; i > §0§; --i )             Ü\itshape\color{gray}/* loopbound min=100 max=100 */\label{lst:lb1}Ü
  [...]

j = §0§;
do {                                    Ü\itshape\color{gray}/* loopbound min=42 max=1234 */\label{lst:lb2}Ü
  [...]
  ++j;
} while ( j < n );
\end{lstlisting}

\begin{lstlisting}[style=code,float=!b,caption={Public Interface of Class \texttt{WIR\_LoopBound}},label=lst:loopbound]
// Default constructor.
WIR_LoopBound::WIR_LoopBound@( int@ min, @int@ max, const WIR_BasicBlock æ&æb,                                                                WIR_LoopBound::LoopControlType ctrl æ=æ LoopControlTypeæ::æunknown,                                           @bool@ chk æ=æ true );Ü\label{lst:loopbound1}Ü

// Copy constructor.
WIR_LoopBound::WIR_LoopBound@(@ const WIR_LoopBound æ&ælb @)@;Ü\label{lst:loopbound2}Ü

// Set the minimum number of iterations of a loop.
@void@ WIR_LoopBound::setMin@( int@ min @)@;Ü\label{lst:loopbound3}Ü

// Get the minimum number of iterations of a loop.
@int@ WIR_LoopBound::getMin@( void )@ const;Ü\label{lst:loopbound4}Ü

// Set the maximum number of iterations of a loop.
@void@ WIR_LoopBound::setMax@( int@ max @)@;Ü\label{lst:loopbound5}Ü

// Get the maximum number of iterations of a loop.
@int@ WIR_LoopBound::getMax@( void )@ const;Ü\label{lst:loopbound6}Ü

// Set the loop that this loop bound is referring to.
@void@ WIR_LoopBound::setLoop@(@ const WIR_BasicBlock æ&æb, LoopControlType ctrl æ=æ LoopControlTypeæ::æunknown,                                 @bool@ chk æ=æ true );Ü\label{lst:loopbound7}Ü

// Get the entry block of the loop that this loop bound is referring to.
const WIR_BasicBlock æ&æWIR_LoopBound::getLoop@( void )@ const;Ü\label{lst:loopbound8}Ü

// Get the control type of the annotated loop.
WIR_LoopBound::LoopControlType WIR_LoopBound::getLoopControlType@( void @) const;Ü\label{lst:loopbound9}Ü

// Enum LoopControlType represents the type of a loop.
enum class WIR_LoopBound::LoopControlType : @char@Ü\label{lst:loopbound10}Ü
{
  headcontrolled,     // A head-controlled loop (for and while-do loops).
  tailcontrolled,     // A tail-controlled loop (do-while loops).
  unknown             // Unknown loop control.
};
\end{lstlisting}

\cref{lst:loopbound} shows the public interface of class \texttt{WIR\_LoopBound}.
Its constructor (cf. \cref{lst:loopbound1}) gets the values for the minimum and maximum iteration counts as first arguments.
Next, the \gls{WIR} basic block \texttt{b} that is the unique entry of the annotated loop according to the structural control flow analysis from \cref{sec:structcflow} needs to be given.
The class of the annotated loop, i.e., whether it is head- or tail-controlled, can be specified using enumeration type \texttt{LoopControlType} (\cref{lst:loopbound10}).
Finally, the Boolean argument \texttt{chk} defaulting to \texttt{true} activates or deactivates a sanity check carried out by the constructor.
This sanity check applies structural analysis and verifies that the given basic block \texttt{b} definitely is entry point of a loop.
If this is not the case, the constructor fails with a runtime error.
There is nothing special to say about the copy constructor from \cref{lst:loopbound2}.

Methods \texttt{setMin} and \texttt{getMin} (\cref{lst:loopbound3,lst:loopbound4}) set and retrieve the minimum iteration count of a loop, resp., the equivalent methods \texttt{setMax} and \texttt{getMax} (\cref{lst:loopbound5,lst:loopbound6}) exist for the maximum iteration count.
The loop entry block can be specified using \texttt{setLoop} (\cref{lst:loopbound7}), the arguments related to the loop's control type and the sanity check are identical to those of the default constructor from \cref{lst:loopbound1}.
The loop's entry block can be retrieved via method \texttt{getLoop} (\cref{lst:loopbound8}), and likewise for the loop control type using \texttt{getLoopControlType} from \cref{lst:loopbound9}.

Flow restrictions are the most powerful and general class of flow facts.
They allow to specify linear inequations relating execution frequencies of different locations in the code with each other.
A typical use case is the annotation of triangular loops where the number of iterations of the inner \texttt{j}-loop below depends on an outer \texttt{i}-loop:

\begin{lstlisting}[style=code,float=!h]
outside: i = §0§;
         while ( i < §10§ ) {
           j = i;
           while ( j < §10§ ) {
             [...]
inside:      ++j;
           }
           ++i;
         }                                    Ü\itshape\color{gray}/* flowrestriction 1*inside <= 55*outside */\label{lst:fr1}Ü
\end{lstlisting}

The flow restriction from \cref{lst:fr1} upper-bounds the execution frequency at label \texttt{inside} by 55 times the execution frequency at label \texttt{outside}.
As can be seen, flow restrictions do a less-equal comparison of terms on the left-hand side with terms on the right-hand side.
These terms consist of a positive integer constant \textit{C} that is multiplied by the supposed execution frequency \textit{freq} at some code location.
In general, there can be many such terms on both sides, and they are all summed up on their respective sides of the inequation:
$\sum\limits_\textit{i} \textit{C}_\textit{i} \ast \textit{freq}_\textit{i} \leq \sum\limits_\textit{j} \textit{C}_\textit{j} \ast \textit{freq}_\textit{j}$.

The code locations to which a flow restriction's summands refer usually reside inside and outside of some loop (nest) as in the above example.
However, these locations can also be spread over different functions, which allows to annotate recursive functions as shown in the following:

\begin{lstlisting}[style=code,float=!h]
         @int@ fib( @int@ n )
         {
           if ( ( n == §0§ ) || ( n == §1§ ) )
             return( §1§ );
           else
             return( fib( n - §1§ ) + fib( n - §2§ ) );
         }

         @int@ main()
         {
           [...]

call:      res = fib( §10§ );                   Ü\itshape\color{gray}/* flowrestriction 1*fib <= 177*call */\label{lst:fr2}Ü

           [...]
         }
\end{lstlisting}

Here, the flow restriction from \cref{lst:fr2} specifies an upper limit of the execution frequency of function \texttt{fib} which is 177 times the execution frequency at label \texttt{call}.

\begin{lstlisting}[style=code,float=!b,caption={Public Interface of Class \texttt{WIR\_FlowRestriction}},label=lst:flowrestr]
// Default constructor with copy-added left- and right-hand sides.
WIR_FlowRestriction::WIR_FlowRestriction@(@ const ästd::listäæ<æ                                                                                                  ästd::pairäæ<æ                                                                                                  @int@,                                                                                                      ästd::reference_wrapperäæ<æconst WIR_BasicBlockæ>>> &ælhs,                                            const ästd::listäæ<æ                                                                                                  ästd::pairäæ<æ                                                                                                  @int@,                                                                                                      ästd::reference_wrapperäæ<æconst WIR_BasicBlockæ>>> &ærhs @)@;Ü\label{lst:flowrestr1}Ü

// Copy constructor.
WIR_FlowRestriction::WIR_FlowRestriction@(@ const WIR_FlowRestriction æ&æfr @)@;Ü\label{lst:flowrestr2}Ü

// Add a summand to a flow restriction's left-hand side.
@void@ WIR_FlowRestriction::addToLHS@( int@ c, const WIR_BasicBlock æ&æb @)@;Ü\label{lst:flowrestr3}Ü

// Add a summand to a flow restriction's right-hand side.
@void@ WIR_FlowRestriction::addToRHS@( int@ c, const WIR_BasicBlock æ&æb @)@;Ü\label{lst:flowrestr4}Ü

// Remove a summand from a flow restriction.
@int@ WIR_FlowRestriction::eraseSummand@(@ const WIR_BasicBlock æ&æb @)@;Ü\label{lst:flowrestr5}Ü

// Replace a summand in a flow restriction.
@void@ WIR_FlowRestriction::replaceSummand@(@ const WIR_BasicBlock æ&æbOld, const WIR_BasicBlock æ&æbNew @)@;Ü\label{lst:flowrestr6}Ü

// Get the left-hand side of a flow restriction.
const ästd::listäæ<æästd::pairäæ<æ@int@, ästd::reference_wrapperäæ<æconst WIR_BasicBlockæ>>> &ægetLHS@( void )@ const;Ü\label{lst:flowrestr7}Ü

// Get the right-hand side of a flow restriction.
const ästd::listäæ<æästd::pairäæ<æ@int@, ästd::reference_wrapperäæ<æconst WIR_BasicBlockæ>>> &ægetRHS@( void )@ const;Ü\label{lst:flowrestr8}Ü

// Determine whether a basic block belongs to a flow restriction's left-hand side.
@bool@ WIR_FlowRestriction::isPartOfLHS@(@ const WIR_BasicBlock æ&æb @)@ const;Ü\label{lst:flowrestr9}Ü

// Determine whether a basic block belongs to a flow restriction's right-hand side.
@bool@ WIR_FlowRestriction::isPartOfRHS@(@ const WIR_BasicBlock æ&æb @)@ const;Ü\label{lst:flowrestr10}Ü

// Determine whether a basic block belongs to a flow restriction at all.
@bool@ WIR_FlowRestriction::isPartOfFlowFact@(@ const WIR_BasicBlock æ&æb @)@ const;Ü\label{lst:flowrestr11}Ü

// Determine whether two flow restrictions are equal.
@bool@ WIR_FlowRestriction::isEqual@(@ const WIR_FlowRestriction æ&æfr @)@ const;Ü\label{lst:flowrestr12}Ü
\end{lstlisting}

\cref{lst:flowrestr} shows the public \gls{API} of class \texttt{WIR\_FlowRestriction}.
Individual summands are represented by pairs of \texttt{int} numbers and (wrapped) references to \gls{WIR} basic blocks that denote the execution frequency at the block's location.
Lists of such pairs finally represent an inequation's left- or right-hand side, resp.
Thus, the constructor from \cref{lst:flowrestr1} obtains two such lists as arguments and builds the flow restriction's left- and right-hand sides from them.
It is enforced that only positive constant factors \textit{C} can be used in both sides of the inequation.
As usual, a copy constructor is provided in \cref{lst:flowrestr2}.

Adding summands to the left-hand side can be done using method \texttt{addToLHS} (\cref{lst:flowrestr3}).
It checks whether a summand referring to the same basic block \texttt{b} already exists on the left-hand side.
If so, its current constant factor is incremented by the value provided in argument \texttt{c}.
Otherwise, a completely new summand will be added.
If the resulting constant factor of the summand becomes negative, the summand is moved from the left- to the right-hand side so that always positive constants show up in a flow restriction.
If the referred basic block \texttt{b}, however, already shows up on the right-hand side, its current constant factor is decremented by \texttt{c} and possibly moved onto the left-hand side if it becomes negative.
Method \texttt{addToRHS} (\cref{lst:flowrestr4}) behaves the same, just with opposite right- and left-hand sides.
Summands whose constant factor becomes zero during \texttt{addToLHS} or \texttt{addToRHS} are erased from the flow restriction, as they are no longer significant.
Method \texttt{eraseSummand} (\cref{lst:flowrestr5}) can be used to remove a summand referring to a given basic block from a flow restriction, regardless of whether the summand is part of the left- or right-hand side.
The summand's current constant factor is returned by this method.
Replacing an old basic block that some summand referred to by a new basic block can be achieved using \texttt{replaceSummand} (\cref{lst:flowrestr6}).
During such a replacement, the previous constant factor of the modified summand remains unchanged.

Access to the full left- and right-hand sides of a flow restriction is provided by methods \texttt{getLHS} and \texttt{getRHS} (\cref{lst:flowrestr7,lst:flowrestr8}), resp.
Whether or not a certain basic block is referred to by some summand on the left- or right-hand side, or any of these, can be checked using \texttt{isPartOfLHS} (\cref{lst:flowrestr9}), \texttt{isPartOfRHS} (\cref{lst:flowrestr10}) and \texttt{isPartOfFlowFact} (\cref{lst:flowrestr11}), resp.
Finally, method \texttt{isEqual} (\cref{lst:flowrestr12}) compares two flow restrictions for strict equality which means that all summands of both flow restrictions must be at exactly the same position in either their left- or right-hand sides, and the summands all must have the same constant factor and must refer to the very same basic block.

This quick overview about the different flow fact classes shows that they all refer to certain positions within \gls{WIR} code -- entry points refer to the beginning of a function, i.e., to a function's very first basic block; loop bounds refer to the entry block of a loop; and a flow restriction refers to an arbitrary number of basic blocks.
This property motivates that solely accessing flow facts from their \gls{WIR} system as shown in \cref{lst:sysffacts} is not always ideal.
Instead, an additional mechanism to access all flow facts referring to a certain basic block might be convenient.
Such a mechanism is provided by \gls{WIR} container class \texttt{WIR\_FlowFactRef} which automatically gets attached to basic blocks.
Whenever a flow fact is inserted into some \gls{WIR} system, it is ensured that appropriate containers are created if they do not yet exist for the referred basic block(s), and that the newly inserted flow fact is introduced to the corresponding container(s).
Upon a removal of a flow fact from a system, it furthermore also gets removed from all \texttt{WIR\_FlowFactRef} containers.

\begin{lstlisting}[style=code,float=t,caption={Flow Fact Handling of Container Class \texttt{WIR\_FlowFactRef}},label=lst:ffref]
// Statically create and get a FlowFactRef container for a given basic block.
static WIR_FlowFactRef æ&æWIR_FlowFactRef::get@(@ WIR_BasicBlock æ&æb @)@;Ü\label{lst:ffref1}Ü

// Get the list of all flow facts referring to the basic block to which this container is attached.
const ästd::listäæ<æästd::reference_wrapperäæ<æWIR_FlowFactæ>> &æWIR_FlowFactRef::getFlowFacts@( void )@ const;Ü\label{lst:ffref2}Ü

// Get the list of all flow facts of a specified type.
templateæ<ætypename Tæ>æ const ästd::listäæ<æästd::reference_wrapperäæ<æTæ>>æ WIR_FlowFactRef::getFlowFacts@( void )@ const;Ü\label{lst:ffref3}Ü
\end{lstlisting}

The public interface of this container class is shown in \cref{lst:ffref}.
It does not feature any public constructors, as the management of these containers is exclusively delegated to \gls{WIR} systems and the individual flow facts themselves.
However, the static \texttt{get} method from \cref{lst:ffref1} should be used in order to obtain a \texttt{WIR\_FlowFactRef} instance for a given basic block.
If no such container is attached to the specified basic block yet, an empty container is created and returned.
Access to all flow facts that somehow refer to that basic block to which a \texttt{WIR\_FlowFactRef} container is attached is given by method \texttt{getFlowFacts} (\cref{lst:ffref2}).
The template-based variant of this method from \cref{lst:ffref3} can be used if only a certain type of flow facts like, e.g., loop bounds, referring to the container's basic block shall be inspected.


%-------------------------------------------------------------------------
\section{Reaching Definitions Analysis}
\label{sec:reachdef}
%-------------------------------------------------------------------------

\begin{tabular}{|l|l|l|l|l|}
  \hline
  \textbf{Analysis Class} & \textbf{Analysis Level} & \textbf{Container Class} & \textbf{Fixing Points} & \textbf{Architectures} \\
  \hhline{|=|=|=|=|=|}
  \texttt{\small WIR\_ReachingDefinitions-} & Functions & \texttt{\small WIR\_Reaching-} & Instructions & Generic, all \\
  \texttt{\small Analysis} & & \texttt{\small Definitions} & & \\
  \hline
\end{tabular}

The first \gls{WIR} data flow analysis presented in this chapter is the reaching definitions analysis realized by class \texttt{WIR\_ReachingDefinitionsAnalysis} (see file \texttt{WIR/analysis/reachingdefinitions/wirreachingdefi\-nitionsanalysis.h}).
For each instruction \texttt{i}, it determines the set of definitions, i.e., of writes to some register \texttt{r}, that may reach the instruction \texttt{i} without being overwritten by some other re-definition in between.
\gls{WIR}'s reaching definitions analysis closely follows the concepts from~\cite[chapter~17.2]{Appel04} and~\cite[chapter~8.1]{Much97}.

Reaching definitions are computed at the scope of individual \gls{WIR} functions \texttt{f}.
For each individual register \texttt{r}, set \textit{defs}( \texttt{r} ) contains all those parameters within the entire function \texttt{f} that somehow modify \texttt{r}, i.e., that either define or define-use \texttt{r} (cf. \cref{ssec:regparams}):
\begin{equation}\label{eq:reachdef1}
  \begin{split}
    \textit{defs}(\ \texttt{r}\ ) := \{\ \text{register parameters}\ \texttt{rp} \in \texttt{f}\ \vert\ & (\texttt{rp.getRegister()} = \texttt{r})\ \wedge \\
                                                                                                        & (\texttt{rp.isDefined()} \vee \texttt{rp.isDefUsed()})\ \}
  \end{split}
\end{equation}
For some instruction \texttt{i}, set \textit{defs}( \texttt{i} ) contains all the registers that are modified by \texttt{i}:
\begin{equation}\label{eq:reachdef2}
  \begin{split}
    \textit{defs}(\ \texttt{i}\ ) := \{\ \text{registers}\ \texttt{r} \in \texttt{f}\ \vert\ \exists \text{register parameter}\ \texttt{rp} \in \texttt{i}:\ & (\texttt{rp.getRegister()} = \texttt{r})\ \wedge \\
                                                                                                                                                             & (\texttt{rp.isDefined()} \vee \texttt{rp.isDefUsed()})\ \}
  \end{split}
\end{equation}
Note the difference between the two sets: the former contains locations in \gls{WIR} code where a register is overwritten, while the latter contains all registers that somehow get overwritten by an instruction.

An instruction \texttt{i} generates new definitions, i.e., by the end of \texttt{i}, the parameters in set \textit{gen}( \texttt{i} ) contain newly written values:
\begin{equation}\label{eq:reachdef3}
  \textit{gen}(\ \texttt{i}\ ) := \{\ \text{register parameters}\ \texttt{rp} \in \texttt{i}\ \vert\ \texttt{rp.isDefined()} \vee \texttt{rp.isDefUsed()}\ \}
\end{equation}
An instruction not only generates new definitions, but it also kills previous definitions.
If a definition of some register \texttt{r} is assumed to reach the beginning of \texttt{i}, but \texttt{i} overwrites this register with a new value, then this previous definition never reaches the end of \texttt{i}. \texttt{i} thus kills the previous definition of \texttt{r}.
All potential definitions of some register \texttt{r} are captured in set \textit{defs}( \texttt{r} ) from \cref{eq:reachdef1}.
Subtracting the set \textit{gen}( \texttt{i} ) of newly generated definitions yields set \textit{kill}( \texttt{i} ):
\begin{equation}\label{eq:reachdef4}
  \textit{kill}(\ \texttt{i}\ ) := \bigcup\limits_{\text{registers}\ \texttt{r}\,\in\,\textit{defs}(\ \texttt{i}\ )} \textit{defs}(\ \texttt{r}\ ) - \textit{gen}(\ \texttt{i}\ )
\end{equation}

Based on the sets \textit{gen} and \textit{kill} per instruction \texttt{i}, the sets of definitions that reach the beginning and the end of \texttt{i}, resp., are defined by the following two data flow equations:
\begin{align}
  \textit{ReachDefsIn}(\ \texttt{i}\ ) :=\ & \bigcup\limits_{\texttt{p}\,\in\,\textit{pred}(\ \texttt{i}\ )} \textit{ReachDefsOut}(\ \texttt{p}\ )\label{eq:reachdef5} \\
  \textit{ReachDefsOut}(\ \texttt{i}\ ) :=\ & \textit{gen}(\ \texttt{i}\ ) \cup (\textit{ReachDefsIn}(\ \texttt{i}\ ) - \textit{kill}(\ \texttt{i}\ ))\label{eq:reachdef6}
\end{align}
\cref{eq:reachdef6} states that everything that is generated by \texttt{i} reaches the end of \texttt{i}, as well as everything that reaches the beginning of \texttt{i} but is not killed by \texttt{i}.
Whatever reaches the end of some predecessors \texttt{p} of \texttt{i} also reaches the beginning of \texttt{i}, as stated by \cref{eq:reachdef5}.

These data flow equations are solved by iteration over a function's \gls{CFG} in reverse postorder within class \texttt{WIR\_\-ReachingDefinitionsAnalysis}.
Since the mechanisms within \gls{WIR} to model \gls{CFG} predecessors and definitions of registers are generic and processor-independent, reachability analysis does not require any adaptations for specific architectures.
Besides the default analysis \gls{API} from \cref{lst:wiranalysis}, the two methods shown in \cref{lst:reacha} are provided.
They control whether the analysis only considers definitions of virtual registers, or whether both virtual and physical registers as well as their precolorings are analyzed.
For this purpose, method \texttt{setVirtualRegistersOnly} (cf. \cref{lst:reacha1}) sets this analysis behavior using the Boolean argument defaulting to \texttt{true}.
The current status of the analysis can be retrieved using \texttt{getVirtualRegistersOnly} (\cref{lst:reacha2}).
In general, reaching definitions analysis is inherently able to deal with arbitrary mixtures of virtual and physical registers, as well as with arbitrary hierarchies of registers (cf. \cref{sssec:regs} on \cpageref{sssec:regs}).
Furthermore, registers defined outside some function \texttt{f} but which serve as inputs to \texttt{f} (cf. \cref{lst:fctargs} on \cpageref{lst:fctargs}) are also properly considered during this analysis.

\begin{lstlisting}[style=code,float=t,caption={Public Interface of Class \texttt{WIR\_ReachingDefinitionsAnalysis}},label=lst:reacha]
// Set whether reaching definitions of virtual registers only are analyzed.
@void@ WIR_ReachingDefinitionsAnalysis::setVirtualRegistersOnly@( bool@ vRegsOnly æ=æ true @)@;Ü\label{lst:reacha1}Ü

// Determine whether reaching definitions of virtual registers only are analyzed.
@bool@ WIR_ReachingDefinitionsAnalysis::getVirtualRegistersOnly@( void )@ const;Ü\label{lst:reacha2}Ü
\end{lstlisting}

The analysis stores the determined sets \textit{ReachDefsIn} from \cref{eq:reachdef5} in \texttt{WIR\_ReachingDefinitions} containers that are attached to \gls{WIR} instructions.
The public interface of this container class is shown in \cref{lst:reach}.
Found reaching definitions can be added using method \texttt{insertReachingDefinition} (cf. \cref{lst:reach1}).
This method does not do any sanity check whether the given register parameter actually is a definition or not, it should thus be used with extreme care.
All identified incoming reaching definitions can be obtained via method \texttt{getReachingDefinitions} (\cref{lst:reach2}).
Similar \texttt{insert} and \texttt{get} methods furthermore exist for incoming reaching function inputs, see \cref{lst:reach3,lst:reach4}.
The sets of register parameters or physical registers returned by \texttt{getReachingDefinitions} and \texttt{getReachingInputs}, resp., are declared along the lines described earlier in \cref{chap:setsmaps}.

\begin{lstlisting}[style=code,float=ht,caption={Reaching Definitions and Inputs Handling of Container Class \texttt{WIR\_ReachingDefinitions}},label=lst:reach]
// Add a new definition to the set of incoming reaching definitions.
@void@ WIR_ReachingDefinitions::insertReachingDefinition@(@ const WIR_RegisterParameter æ&ærp @)@;Ü\label{lst:reach1}Ü

// Get the set of incoming reaching definitions.
const WIR_RegisterParameterSet æ&æWIR_ReachingDefinitions::getReachingDefinitions@( void )@;Ü\label{lst:reach2}Ü

// Add a new input register to the set of incoming reaching function inputs.
@void@ WIR_ReachingDefinitions::insertReachingInput@(@ const WIR_PhysicalRegister æ&æp @)@;Ü\label{lst:reach3}Ü

// Get the set of incoming reaching function inputs.
const WIR_PhysicalRegisterSet æ&æWIR_ReachingDefinitions::getReachingInputs@( void )@;Ü\label{lst:reach4}Ü
\end{lstlisting}


%-------------------------------------------------------------------------
\section{Available Definitions Analysis}
\label{sec:availdef}
%-------------------------------------------------------------------------

\begin{tabular}{|l|l|l|l|l|}
  \hline
  \textbf{Analysis Class} & \textbf{Analysis Level} & \textbf{Container Class} & \textbf{Fixing Points} & \textbf{Architectures} \\
  \hhline{|=|=|=|=|=|}
  \texttt{\small WIR\_AvailableDefinitions-} & Functions & \texttt{\small WIR\_Available-} & Instructions & Generic, all \\
  \texttt{\small Analysis} & & \texttt{\small Definitions} & & \\
  \hline
\end{tabular}

Available definitions (cf. file \texttt{WIR/analysis/availabledefinitions/wiravailabledefinitions.h}) analysis is very similar to reaching definitions analysis.
While reaching definitions is a may analysis, available definitions analysis answers the question which register definitions definitely must reach an instruction \texttt{i} along all control flow paths~\cite[chapter~17.2]{Appel04}.

The previous definitions of the \textit{gen} and \textit{kill} sets from \cref{eq:reachdef3,eq:reachdef4} remain unchanged, only the data flow equations that propagate analysis results through a function's \gls{CFG} change slightly:
\begin{align}
  \textit{AvailDefsIn}(\ \texttt{i}\ ) :=\ & \bigcap\limits_{\texttt{p}\,\in\,\textit{pred}(\ \texttt{i}\ )} \textit{AvailDefsOut}(\ \texttt{p}\ )\label{eq:availdef1} \\
  \textit{AvailDefsOut}(\ \texttt{i}\ ) :=\ & \textit{gen}(\ \texttt{i}\ ) \cup (\textit{AvailDefsIn}(\ \texttt{i}\ ) - \textit{kill}(\ \texttt{i}\ ))\label{eq:availdef2}
\end{align}
As can be seen from \cref{eq:availdef1}, available definitions analysis applies the intersection to all predecessor's outgoing available definitions, which is the natural formalization for a must analysis.
The previous use of the union operator in \cref{eq:reachdef5}, however, is the distinguishing feature of a may analysis.
Other than that, \gls{WIR} class \texttt{WIR\_AvailableDefinitionsAnalysis} behaves identical to what has been discussed previously in \cref{sec:reachdef}.
In analogy to \cref{lst:reacha}, the available definitions analysis can also be configured to consider virtual register definitions only.

For the sake of completeness, \cref{lst:avail} shows the public \gls{API} of the analysis' container class \texttt{WIR\_Avail\-able\-Definitions}, which gets attached to \gls{WIR} instructions again.
There is no difference at all compared to \cref{lst:reach} so that we omit the presentation of the container's methods here.

\begin{lstlisting}[style=code,float=ht,caption={Available Definitions and Inputs Handling of Container Class \texttt{WIR\_AvailableDefinitions}},label=lst:avail]
// Add a new definition to the set of incoming available definitions.
@void@ WIR_AvailableDefinitions::insertAvailableDefinition@(@ const WIR_RegisterParameter æ&ærp @)@;Ü\label{lst:avail1}Ü

// Get the set of incoming available definitions.
const WIR_RegisterParameterSet æ&æWIR_AvailableDefinitions::getAvailableDefinitions@( void )@;Ü\label{lst:avail2}Ü

// Add a new input register to the set of incoming available function inputs.
@void@ WIR_AvailableDefinitions::insertAvailableInput@(@ const WIR_PhysicalRegister æ&æp @)@;Ü\label{lst:avail3}Ü

// Get the set of incoming available function inputs.
const WIR_PhysicalRegisterSet æ&æWIR_AvailableDefinitions::getAvailableInputs@( void )@;Ü\label{lst:avail4}Ü
\end{lstlisting}


%-------------------------------------------------------------------------
\section{Def-Use/Use-Def Chains Analysis}
\label{sec:duudchains}
%-------------------------------------------------------------------------

\begin{tabular}{|l|l|l|l|l|}
  \hline
  \textbf{Analysis Class} & \textbf{Analysis Level} & \textbf{Container Class} & \textbf{Fixing Points} & \textbf{Architectures} \\
  \hhline{|=|=|=|=|=|}
  \texttt{\small WIR\_DUUDChainAnalysis} & Functions & \texttt{\small WIR\_DUUDChain} & Register Parameters & Generic, all \\
  \hline
\end{tabular}

The combined def-use and use-def chain analysis of \gls{WIR} realized by class \texttt{WIR\_DUUDChainAnalysis} (see file \texttt{WIR/analyses/wirduudchainanalysis.h}) bases on the reaching definitions analysis from \cref{sec:reachdef}.
While reaching definitions provides the set of all register definitions that may reach the beginning of an instruction, def-use and use-def chains represent the same information but limited to only one given register parameter \texttt{rp}~\cite[chapter~17.4]{Appel04} and~\cite[chapter~8.10]{Much97}.
More precisely, a def-use chain is defined as follows:

\begin{definition}[Def-Use Chain] \label{def:duchain}
  For a register parameter \texttt{\emph{rp}}$_\textit{def}$ defining some register \texttt{\emph r}, its \emph{def-use chain} is the set of all register parameters \texttt{\emph{rp}}$_\textit{use}$ using the very same register \texttt{\emph r} such that the definition from \texttt{\emph{rp}}$_\textit{def}$ may reach \texttt{\emph{rp}}$_\textit{use}$.
  \begin{align}
    \forall\text{\emph{register parameters}\ }\texttt{\emph{rp}}_\textit{def}&:\ \texttt{\emph{rp}}_\textit{def}\texttt{\emph{.getUsage()}} \in \{ \texttt{\emph{def}},\ \texttt{\emph{defuse}}\}:\nonumber \\
    \textit{DU-Chain}(\ \texttt{\emph{rp}}_\textit{def}\ ) &:= \{ \text{\emph{register parameters}\ }\texttt{\emph{rp}}_\textit{use}\ \vert\ (\texttt{\emph{rp}}_\textit{use}\texttt{\emph{.getUsage()}} \in \{ \texttt{\emph{use}},\ \texttt{\emph{defuse}}\})\ \wedge\label{eq:duchain} \\
    & \ \ \ \ \ \ \ \ \ \ \ \ \ (\texttt{\emph{rp}}_\textit{def}\texttt{\emph{.getRegister()}} = \texttt{\emph{rp}}_\textit{use}\texttt{\emph{.getRegister()}})\ \wedge\nonumber \\
    & \ \ \ \ \ \ \ \ \ \ \ \ \ (\texttt{\emph{rp}}_\textit{def} \in \text{ReachDefsIn}(\ \texttt{\emph{rp}}_\textit{use}\texttt{\emph{.getOperation().getInstruction()}}\ ))\}\nonumber
  \end{align}
\end{definition}

Similarly, a use-def chain for a used register parameter consists of all definitions of the very same register that may reach the use:

\begin{definition}[Use-Def Chain] \label{def:udchain}
  For a register parameter \texttt{\emph{rp}}$_\textit{use}$ using some register \texttt{\emph r}, its \emph{use-def chain} is the set of all register parameters \texttt{\emph{rp}}$_\textit{def}$ defining the very same register \texttt{\emph r} such that the definitions from \texttt{\emph{rp}}$_\textit{def}$ may reach \texttt{\emph{rp}}$_\textit{use}$.
  \begin{align}
    \forall\text{\emph{register parameters}\ }\texttt{\emph{rp}}_\textit{use}&:\ \texttt{\emph{rp}}_\textit{use}\texttt{\emph{.getUsage()}} \in \{ \texttt{\emph{use}},\ \texttt{\emph{defuse}}\}:\nonumber \\
    \textit{UD-Chain}(\ \texttt{\emph{rp}}_\textit{use}\ ) &:= \{ \text{\emph{register parameters}\ }\texttt{\emph{rp}}_\textit{def}\ \vert\ (\texttt{\emph{rp}}_\textit{def}\texttt{\emph{.getUsage()}} \in \{ \texttt{\emph{def}},\ \texttt{\emph{defuse}}\})\ \wedge\label{eq:udchain} \\
    & \ \ \ \ \ \ \ \ \ \ \ \ \ (\texttt{\emph{rp}}_\textit{use}\texttt{\emph{.getRegister()}} = \texttt{\emph{rp}}_\textit{def}\texttt{\emph{.getRegister()}})\ \wedge\nonumber \\
    & \ \ \ \ \ \ \ \ \ \ \ \ \ (\texttt{\emph{rp}}_\textit{def} \in \text{ReachDefsIn}(\ \texttt{\emph{rp}}_\textit{use}\texttt{\emph{.getOperation().getInstruction()}}\ ))\}\nonumber
  \end{align}
\end{definition}

The def-use/use-def chain analysis thus first performs a reaching definitions analysis and then restructures its results such that the chains from \cref{def:duchain,def:udchain} are formed.
Of course, hierarchies and precolorings of registers are properly considered again, as well as registers serving as function input.
Not surprisingly, def-use/use-def chain analysis is also fully generic and does not require any processor-specific adaptations.
Due to its strong relationship to reaching definitions, the public interface of class \texttt{WIR\_DUUDChainAnalysis} is identical to that already shown earlier in \cref{lst:reacha}.

\begin{lstlisting}[style=code,float=!b,caption={Def-Use/Use-Def Chain and Inputs Handling of Container Class \texttt{WIR\_DUUDChain}},label=lst:duudchain]
// Add a new use to the def-use chain.
@void@ WIR_DUUDChain::insertDUChain@(@ const WIR_RegisterParameter æ&ærp @)@;Ü\label{lst:duudchain1}Ü

// Get the def-use chain.
const WIR_RegisterParameterSet æ&æWIR_DUUDChain::getDUChains@( void )@;Ü\label{lst:duudchain2}Ü

// Add a new definition to the use-def chain.
@void@ WIR_DUUDChain::insertUDChain@(@ const WIR_RegisterParameter æ&ærp @)@;Ü\label{lst:duudchain3}Ü

// Get the use-def chain.
const WIR_RegisterParameterSet æ&æWIR_DUUDChain::getUDChains@( void )@;Ü\label{lst:duudchain4}Ü

// Add a new defined input register to the use-def chain.
@void@ WIR_DUUDChain::insertUDInput@(@ const WIR_PhysicalRegister æ&æp @)@;Ü\label{lst:duudchain5}Ü

// Get the use-def chain for external function inputs.
const WIR_PhysicalRegisterSet æ&æWIR_DUUDChain::getUDInputs@( void )@;Ü\label{lst:duudchain6}Ü
\end{lstlisting}

Since def-use and use-def chains are defined for individual register parameters, the analysis attaches containers of class \texttt{WIR\_DUUDChain} to all such register parameters within an analyzed \gls{WIR} function.
\cref{lst:duudchain} shows the public interface of this container class.
The class provides the method \texttt{insertDUChain} to add a use to the def-use chain of the register parameter to which the container is attached (cf. \cref{lst:duudchain1}), while method \texttt{getDUChains} (\cref{lst:duudchain2}) retrieves the register parameter set forming a def-use chain.
Similarly, methods \texttt{insertUDChain} (\cref{lst:duudchain3}) and \texttt{getUDChain} (\cref{lst:duudchain4}) can be used to handle use-def chains.
Since use-def chains also need to consider external function inputs provided by dedicated physical registers, methods \texttt{insertUDInput} (\cref{lst:duudchain5}) and \texttt{getUDInputs} (\cref{lst:duudchain6}) exist to extend and retrieve these sets of input registers belonging to a use-def chain.


%-------------------------------------------------------------------------
\section{Liveness Analysis}
\label{sec:liveness}
%-------------------------------------------------------------------------

\begin{tabular}{|l|l|l|l|l|}
  \hline
  \textbf{Analysis Class} & \textbf{Analysis Level} & \textbf{Container Class} & \textbf{Fixing Points} & \textbf{Architectures} \\
  \hhline{|=|=|=|=|=|}
  \texttt{\small WIR\_LivenessAnalysis} & Functions & \texttt{\small WIR\_LiveOut} & Instructions & Generic, all \\
  \hline
\end{tabular}

Liveness analysis~\cite[chapter~10]{Appel04} and~\cite[chapter~14.1.3]{Much97} is a prerequisite for any kind of register allocation.
It determines from where to where within an analyzed \gls{WIR} function certain registers are alive, i.e., contain useful values that might be needed as input by some later instruction.
Using this information, a register allocator can determine whether the liveness ranges of two virtual registers overlap somewhere and can then ensure that these two virtual registers are never allocated to the same physical register.

The analysis is realized by class \texttt{WIR\_LivenessAnalysis} which can be found in file \texttt{WIR/analyses/liveness/ wirlivenessanalysis.h}.
Within \gls{WIR}, liveness is considered at the level of individual instructions.
In analogy to reaching definitions analysis (cf. \cref{sec:reachdef}), it bases on the sets of registers that are defined or used by some instruction \texttt{i}:
\begin{align}
  \textit{defs}(\ \texttt{i}\ ) := \{\ \text{registers}\ \texttt{r} \in \texttt{f}\ \vert\ \exists \text{register parameter}\ \texttt{rp} \in \texttt{i}:\ & (\texttt{rp.getRegister()} = \texttt{r})\ \wedge\label{eq:liveness1} \\
                                                                                                                                                            & (\texttt{rp.isDefined()} \vee \texttt{rp.isDefUsed()})\ \} \nonumber \\
  \textit{uses}(\ \texttt{i}\ ) := \{\ \text{registers}\ \texttt{r} \in \texttt{f}\ \vert\ \exists \text{register parameter}\ \texttt{rp} \in \texttt{i}:\ & (\texttt{rp.getRegister()} = \texttt{r})\ \wedge\label{eq:liveness2} \\
                                                                                                                                                            & (\texttt{rp.isUsed()} \vee \texttt{rp.isDefUsed()})\ \} \nonumber
\end{align}
A register \texttt{r} is said to be alive along a \gls{CFG} edge if there is a directed path from that edge to some use of \texttt{r} that contains no other re-definition of \texttt{r}.
The sets of registers that are alive at the beginning and the end of an instruction \texttt{i} are defined by the following two data flow equations:
\begin{align}
  \textit{LiveIn}(\ \texttt{i}\ ) :=\ & \textit{uses}(\ \texttt{i}\ ) \cup (\textit{LiveOut}(\ \texttt{i}\ ) - \textit{defs}(\ \texttt{i}\ ))\label{eq:liveness3} \\
  \textit{LiveOut}(\ \texttt{i}\ ) :=\ & \bigcup\limits_{\texttt{s}\,\in\,\textit{succ}(\ \texttt{i}\ )} \textit{LiveIn}(\ \texttt{s}\ )\label{eq:liveness4}
\end{align}
\cref{eq:liveness3} states that everything that is used by \texttt{i} is live-in, as well as everything that is live-out at \texttt{i} and that is not re-defined by \texttt{i}.
Whatever is live-in at the beginning of some successors \texttt{s} of \texttt{i} is also live-out at \texttt{i} which is modeled by \cref{eq:liveness4}.

These explanations show that liveness gets propagated through a function's \gls{CFG} in the reverse sense of the graph's edges, i.e., liveness analysis is a backward analysis.
These equations are thus solved by iteration through the \gls{CFG} in reverse-topological order within class \texttt{WIR\_LivenessAnalysis}.
Once more, liveness analysis is generic and entirely processor-independent.
The public \gls{API} of the analysis contains no additional methods beyond the default interface from \cref{lst:wiranalysis}.

\begin{lstlisting}[style=code,float=ht,caption={Liveness Handling of Container Class \texttt{WIR\_LiveOut}},label=lst:liveness]
// Add a register to the live-out set.
@void@ WIR_LiveOut::insertRegister@(@ const WIR_BaseRegister æ&ær @)@;Ü\label{lst:liveness1}Ü

// Get the live-out set.
const WIR_RegisterSet æ&æWIR_LiveOut::getRegisters@( void )@ const;Ü\label{lst:liveness2}Ü

// Determine whether the live-out set contains an item with the given ID.
@bool@ WIR_LiveOut::containsRegister@(@ WIR_id_t id @)@ const;Ü\label{lst:liveness3}Ü

// Determine whether the live-out set contains the given register.
@bool@ WIR_LiveOut::containsRegister@(@ const WIR_BaseRegister æ&ær @)@ const;Ü\label{lst:liveness4}Ü

// Find the register with the given ID in the live-out set.
WIR_RegisterSetæ::const_iteratoræ WIR_LiveOut::findRegister@(@ WIR_id_t id @)@ const;Ü\label{lst:liveness5}Ü

// Find the given register in the live-out set.
WIR_RegisterSetæ::const_iteratoræ WIR_LiveOut::findRegister@(@ const WIR_BaseRegister æ&ær @)@ const;Ü\label{lst:liveness6}Ü

// Remove the specified register from the live-out set.
@void@ WIR_LiveOut::eraseRegister@(@ const WIR_BaseRegister æ&ær @)@;Ü\label{lst:liveness7}Ü

// Remove all registers from the live-out set.
@void@ WIR_LiveOut::clearRegisters@( void )@;Ü\label{lst:liveness8}Ü
\end{lstlisting}

The results of liveness analysis are stored in objects of the container class \texttt{WIR\_LiveOut} that are attached to individual \gls{WIR} instructions.
As the name suggests, only the live-out sets are stored for the sake of efficiency.
\cref{lst:liveness} shows the public \gls{API} of this container class, which is kept very similar to the concepts already introduced in \cref{chap:lists}.

Method \texttt{insertRegister} (cf. \cref{lst:liveness1}) can be used to manually add a register to a live-out set.
The complete live-out set can be retrieved using \texttt{getRegisters} (\cref{lst:liveness2}).
The \texttt{containsRegister} methods (\cref{lst:liveness3,lst:liveness4}) and the \texttt{findRegister} methods (\cref{lst:liveness5,lst:liveness6}) can be used to check whether a register is part of a live-out set, and to find it in the set, resp.
Finally, elements can be removed from a live-out set using \texttt{eraseRegister} (\cref{lst:liveness7}) and \texttt{clearRegisters} (\cref{lst:liveness8}), resp.


%-------------------------------------------------------------------------
\section{Bit-True Data Flow Analysis}
\label{sec:bitdflow}
%-------------------------------------------------------------------------

\begin{tabular}{|l|l|l|l|l|}
  \hline
  \textbf{Analysis Class} & \textbf{Analysis Level} & \textbf{Container Class} & \textbf{Fixing Points} & \textbf{Architectures} \\
  \hhline{|=|=|=|=|=|}
  \texttt{\small WIR\_BitDFA} & Functions & \texttt{\small WIR\_BitValues} & Register Parameters, & RV32IMC, \\
  \texttt{\small RV\_BitDFA}, \texttt{\small TC\_BitDFA} &           &                                & Immediate Parameters & TC1.3, TC1.3.1 \\
  \hline
\end{tabular}

While all the data flow analyses presented in the previous \crefrange{sec:reachdef}{sec:liveness} considered complete registers as a whole, it sometimes might be useful to obtain more fine-grained information about the potential contents of individual bits of a register.
This is the purpose of the bit-true data flow analysis~\cite{Wag06} whose generic implementation is provided by the abstract base class \texttt{WIR\_BitDFA} from file \texttt{WIR/analyses/bit/wirbitdfa.h}.
In contrast to all analyses introduced previously in this chapter, bit-true data flow analysis requires processor-specific implementation details.
At the time of writing of this document, ports of the analysis exist for the RISC-V RV32IMC \gls{ISA} (class \texttt{RV32\_BitDFA} from file \texttt{WIR/arch/riscv/analyses/bit/rv32bitdfa.h}) and for both Infineon TriCore models (class \texttt{TC\_BitDFA} from file \texttt{WIR/arch/tricore/analyses/bit/tcbitdfa.h}).

Information about a single bit is represented using the enumeration class \texttt{WIR\_L4} (see file \texttt{WIR/analyses/bit/ wirl4.h}) shown in Listing~\ref{fig:l4enum}.
The elements of this enumeration form a partial order depicted in \cref{fig:l4hasse}.
A completely unknown bit where neither its value nor its origin are known is represented by element \texttt{bU} (cf. \cref{lst:l47}).
If is known that a bit provably contains the value 0 or 1, it is represented by elements \texttt{b0} or \texttt{b1} (cf. \cref{lst:l43,lst:l44}), resp.
A bit whose value is irrelevant so that it does not have any influence on the result computed by some \gls{WIR} code is modeled by \texttt{bX} (\cref{lst:l42}).
If the actual value of a bit cannot be determined, but if it is known that the bit must be identical with another bit at a different location, it is represented by \texttt{bL} (\cref{lst:l45}).
Finally, enumeration element \texttt{bN} (\cref{lst:l46}) models that a bit's value is unknown, but the bit is the negation of another bit at a different location.

\begin{figure}[t]
  \begin{subfigure}[b]{0.75\textwidth}
    \begin{lstlisting}[style=code]
// Enum WIR_L4 represents the L4 partial order.
enum class WIR_L4 : @char@Ü\label{lst:l41}Ü
{
  bX,       // Arbitrary bit value, don't care.Ü\label{lst:l42}Ü
  b0,       // Exactly known bit value of '0'.Ü\label{lst:l43}Ü
  b1,       // Exactly known bit value of '1'.Ü\label{lst:l44}Ü
  bL,       // Unknown bit value but known location of the bit's origin.Ü\label{lst:l45}Ü
  bN,       // Unknown bit value but negated known location of the bit's origin.Ü\label{lst:l46}Ü
  bU        // Unknown bit value and unknown origin.Ü\label{lst:l47}Ü
};
\end{lstlisting}
    \caption{Enumeration Class \texttt{WIR\_L4}}
    \label{fig:l4enum}
  \end{subfigure}
  \begin{subfigure}[b]{0.23\textwidth}
    \begin{center}
      \includegraphics[width=0.84\textwidth]{l4.png}
    \end{center}
    \caption{Hasse Diagram}
    \label{fig:l4hasse}
  \end{subfigure}
  \caption{Partial Order \texttt{L4}}
  \label{fig:l4}
\end{figure}

According to the Hasse diagram shown in \cref{fig:l4hasse}, element \texttt{bU} represents the least degree of information about a bit.
Not knowing a bit's value but being aware of its original location is useful so that \texttt{bL} and \texttt{bN} have a higher position in the Hasse diagram.
Even more information is represented by the next level of the partial order, consisting of the actual bit values \texttt{b0} and \texttt{b1}.
Finally, the largest amount of information about a bit is given if it does not care at all so that \texttt{bX} forms the top in \cref{fig:l4hasse}.
These levels of the partial order represent a \texttt{<} operator along the directed edges of \cref{fig:l4hasse} that is also declared in file \texttt{wirl4.h}.

(Negated) locations of a bit's origin are represented by class \texttt{WIR\_Location} from file \texttt{WIR/analyses/bit/wir\-location.h}.
Similar to classical, register-level data flow analysis, such a location regularly refers to an operation's register parameter that defines the currently considered bit.
But in addition to this, the position of the bit from its original location is furthermore provided.

\begin{wirExample}\label{ex:bit}
Consider the following code snippet based on the RISC-V RV32 processor model introduced in \cref{ssec:rv32model} as example:
\begin{lstlisting}[style=code]
        srli            x_3, x10, §2§                Ü\itshape\color{gray}\# x\_3 := x10 >{}> 2\label{lst:rv32bit1}Ü
        andi            x_4, x_3, §7§                Ü\itshape\color{gray}\# x\_4 := x\_3 \& 7\label{lst:rv32bit2}Ü
\end{lstlisting}
\cref{lst:rv32bit1} takes an external function argument provided in the physical argument register \texttt{x10}, shifts it right by 2 bit positions, and puts the result of this shift in virtual register \texttt{x\_3}.
Hereafter, \cref{lst:rv32bit2} applies a bitwise AND operation by the constant 7 to this virtual register \texttt{x\_3}.
The result of this AND operation is put in virtual register \texttt{x\_4}.

Here and in the following, the bits of a register are enumerated from the least-significant position 0 to the most-significant position 31.
Then, the actual value of bit 0 of register \texttt{x\_4} is unknown, since the bit-true data flow analysis of \gls{WIR} is an intraprocedural analysis that does not track external function arguments.
However, it becomes evident from the above code that this bit 0 of \texttt{x\_4} must be the very same as bit 2 of register \texttt{x10} due to the nature of the \texttt{srli} and \texttt{andi} operations.
Likewise, bit 1 of \texttt{x\_4} is equal to bit 3 of \texttt{x10}, and bit 2 of \texttt{x\_4} matches with bit 4 of \texttt{x10}.

The additional line
\begin{lstlisting}[style=code,firstnumber=3]
        xori            x_5, x_4, §4§                Ü\itshape\color{gray}\# x\_5 := x\_4 \^\ 4\label{lst:rv32bit3}Ü
\end{lstlisting}
has the effect that one bit of \texttt{x\_4} gets inverted.
Thus, bit 2 of \texttt{x\_5} finally is the negation of bit 4 of \texttt{x10}.
\end{wirExample}

\begin{lstlisting}[style=code,float=!b,caption={Location Handling by Class \texttt{WIR\_Location}},label=lst:location]
// Default constructor creating a location referring to a bit produced by a given register parameter.
WIR_Location::WIR_Location@(@ const WIR_RegisterParameter æ&ærp, @unsigned int@ b @)@;Ü\label{lst:location1}Ü

// Default constructor creating a location referring to a bit of a symbol's address.
WIR_Location::WIR_Location@(@ const WIR_Symbol æ&æs, @unsigned int@ b @)@;Ü\label{lst:location2}Ü

// Get a location's type.
WIR_LocationType WIR_Location::getType@( void )@ const;Ü\label{lst:location3}Ü

// Determine whether a location refers to a register parameter.
@bool@ WIR_Location::isRegisterParameter@( void )@ const;Ü\label{lst:location4}Ü

// Determine whether a location refers to a symbol.
@bool@ WIR_Location::isSymbol@( void )@ const;Ü\label{lst:location5}Ü

// Get a location's register parameter.
const WIR_RegisterParameter æ&æWIR_Location::getRegisterParameter@( void )@ const;Ü\label{lst:location6}Ü

// Get a location's symbol.
const WIR_Symbol æ&æWIR_Location::getSymbol@( void )@ const;Ü\label{lst:location7}Ü

// Get a location's bit position.
@unsigned int@ WIR_Location::getBitPosition@( void )@ const;Ü\label{lst:location8}Ü

// Equality operator for locations.
@bool@ WIR_Location::operator == @(@ const WIR_Location æ&æl @)@ const;Ü\label{lst:location9}Ü

// Inequality operator for locations.
@bool@ WIR_Location::operator != @(@ const WIR_Location æ&æl @)@ const;Ü\label{lst:location10}Ü

// Enum WIR_LocationType represents different types of locations.
enum class WIR_LocationType : @char@Ü\label{lst:location11}Ü
{
  reg,          // A location referring to a register parameter.
  sym           // A location referring to a symbol.
};
\end{lstlisting}

The interface of class \texttt{WIR\_Location} is shown in \cref{lst:location}.
A location referring to a certain bit position within some defined register parameter can be created with the constructor shown in \cref{lst:location1}.
A location's register parameter can be retrieved using \texttt{getRegisterParameter} (\cref{lst:location6}), and its bit position is returned by \texttt{getBitPosition} (\cref{lst:location8}).

Locations not necessarily refer to defined register parameters only, but they can also refer to \gls{WIR} symbols, i.e., to a single bit of some memory address.
For this purpose, the constructor from \cref{lst:location2} can be used, and method \texttt{getSymbol} (\cref{lst:location7}) retrieves a location's symbol.

Whether a location refers either to a register parameter or a symbol can be checked using  methods \texttt{getType}, \texttt{isRegisterParameter} or \texttt{isSymbol} (\crefrange{lst:location3}{lst:location5}), resp., where the enumeration type returned by \texttt{getType} is shown in \cref{lst:location11}.
As usual within \gls{WIR}, \texttt{getRegisterParameter} fails with a runtime error if the location refers to a symbol, and vice versa for \texttt{getSymbol}.

Finally, the operators \texttt{==} and \texttt{!=} from \cref{lst:location9,lst:location10} can be used to compare two locations for equality or inequality, resp.
Here, two locations are equal iff they both have the same type and refer to the very same bit of the very same register parameter or symbol, resp.

\begin{figure}[t]
  \begin{center}
    \includegraphics[width=0.8\textwidth]{bitdfg.png}
  \end{center}
  \caption{Example of Bit-True Data Flow Graph}
  \label{fig:bitdfg}
\end{figure}

The concepts regarding the partial order \texttt{L4} and locations refer to one single bit to be analyzed.
As \gls{WIR} operations regularly use complete registers of a certain bit width as parameters (see also the previous example on \cpageref{ex:bit}), single bit values from \texttt{L4} are combined to form values that represent all bits of a register.
These combined values (called down and up values in the following) are associated with the edges of the bit-true \gls{DFG} which is defined as follows:

\begin{definition}[Bit-true \acrlong{DFG}] \label{def:bitdfg}
  A bit-true \emph{data flow graph} $\textit{DFG} = (V, E, \mathit{Dn}, \mathit{Up})$ is a directed graph with
  \begin{itemize}
    \item $V = \{$ $v_1, \dots, v_n$ $\}$ being the set of all nodes, where each node represents either
      \begin{itemize}
        \item an input value (external registers or immediates), or
        \item a single \gls{WIR} operation, or
        \item an output value,
      \end{itemize}
    \item $E = \{$ $(v_i^x \rightarrow v_j^y)$ $\vert$ parameter $x$ of node $i$ defines a value used by parameter $y$ of node $j$ $\}$,
    \item $\mathit{Dn}: E \rightarrow \texttt{\emph{L4}}^k$ and $\mathit{Up}: E \rightarrow \texttt{\emph{L4}}^k$ are the down and up values representing bit-true analysis results per graph edge where the determined values are $k$ bits wide.
  \end{itemize}
\end{definition}

\begin{wirExample}\label{ex:bitdfg}
For the three operations from the previous example, the bit-true \gls{DFG} and all analyzed down and up values are depicted in \cref{fig:bitdfg}.

Here, the immediate value 7 required by the \texttt{andi} operation is represented as ``\texttt{0000 0000 0111}'' which is the straightforward outcome due to the binary representation of the number 7 using a width of 12 bits, as mandated by the RISC-V \gls{ISA}.
The immediate 2 used by \texttt{srli} is also modeled using 12 bits, but the most-significant 7 bits are set to \texttt{bX} here: ``\texttt{XXXX XXX0 0010}''.
This is because the analysis determines that shift amounts in the RISC-V \gls{ISA} can only take values between 0 and 31 so that only the least-significant 5 bits are relevant.
The most-significant bits of shift amounts thus don't care.

Values analyzed for \gls{DFG} edges representing RV32I registers are all 32 bits wide.
The edge from \texttt{andi} to \texttt{xori} along virtual register \texttt{x\_4} shows that only the least-significant 3 bits of \texttt{x\_4} are relevant for the given example, all other most-significant bits are set to \texttt{bX}.
Furthermore, these least-significant 3 bits of \texttt{x\_4} are identical with bits 2-4 of input register \texttt{x10} which is represented by the locations ``\texttt{L(x10[4 - 2])}'' in \cref{fig:bitdfg}.

As already discussed in the previous example, the result finally computed by \texttt{xori} in \texttt{x\_5} contains the original bits 2 and 3 of \texttt{x10} in the least-significant positions, and the inverted bit 4 of \texttt{x10} next: ``\texttt{N(x10[4])L(x10[3 - 2])}''.
\end{wirExample}

\begin{lstlisting}[style=code,float=!b,caption={Constructors of Class \texttt{WIR\_UpDownValue}},label=lst:updnconstruction]
// Default constructor for unsigned up/down values of a given width.
explicit WIR_UpDownValue::WIR_UpDownValue@( unsigned int@ w @)@;Ü\label{lst:updnconstruction1}Ü

// Default constructor for up/down values with all bits set to some L4 value and given width and signedness.
WIR_UpDownValue::WIR_UpDownValue@(@ WIR_L4 v, @unsigned int@ w, @bool@ s æ=æ true @)@;Ü\label{lst:updnconstruction2}Ü

// Default constructor for up/down values initialized to some given immediate value.
explicit WIR_UpDownValue::WIR_UpDownValue@(@ const WIR_BaseImmediateParameter æ&æi @)@;Ü\label{lst:updnconstruction3}Ü

// Default constructor for up/down values initialized to some register location and signedness.
WIR_UpDownValue::WIR_UpDownValue@(@ const WIR_RegisterParameter æ&ærp, @bool@ s æ=æ true, @bool@ n æ=æ false @)@;Ü\label{lst:updnconstruction4}Ü

// Copy constructor.
WIR_UpDownValue::WIR_UpDownValue@(@ const WIR_UpDownValue æ&æv @)@;Ü\label{lst:updnconstruction5}Ü

// Move constructor.
WIR_UpDownValue::WIR_UpDownValue@(@ WIR_UpDownValue æ&&æv @)@;Ü\label{lst:updnconstruction6}Ü
\end{lstlisting}

All down and up values determined during bit-true data flow analysis are modeled by class \texttt{WIR\_UpDownValue} (cf. file \texttt{WIR/analyses/bit/wirupdownvalue.h}).
The constructors of this class are shown in \cref{lst:updnconstruction}.
The most simple constructor from \cref{lst:updnconstruction1} creates an unsigned up/down value of a given bit width whose bits are all initialized with \texttt{bU}.
If some other \texttt{L4} value shall be used for initialization, the variant from \cref{lst:updnconstruction2} can be used that accepts this \texttt{L4} value as argument, as well as its bit width and a Boolean flag defaulting to \texttt{true} that indicates whether the up/down value is interpreted as signed or unsigned value, resp.
Up/down values representing immediate values can conveniently be created with the constructor from \cref{lst:updnconstruction3}.
The actual value as well as its bit width and signedness are directly taken from the given immediate parameter.
The constructor from \cref{lst:updnconstruction4} creates an up/down value that refers to all bits of a given register parameter so that all elements of such an up/down value are initialized with \texttt{bL}.
The second argument of this constructor again denotes the value's signedness, and its last argument denotes whether \texttt{bN} shall be used for initialization.
Finally, the usual copy and move constructors also exist, see \cref{lst:updnconstruction5,lst:updnconstruction6}.

\begin{lstlisting}[style=code,float=!t,caption={Value Handling of Class \texttt{WIR\_UpDownValue}},label=lst:updnvalues]
// Set all bits to a given L4 value different from bL and bN.
@void@ WIR_UpDownValue::setAllBits@(@ WIR_L4 v @)@;Ü\label{lst:updnvalues1}Ü

// Set all bits to a given L4 value bL or bN
@void@ WIR_UpDownValue::setAllBits@(@ WIR_L4 v, const WIR_Location æ&æl @)@;Ü\label{lst:updnvalues2}Ü

// Set a single bit at a given position to an L4 value different from bL and bN.
@void@ WIR_UpDownValue::setBit@( unsigned int@ p, WIR_L4 v @)@;Ü\label{lst:updnvalues3}Ü

// Set a single bit at a given position to an L4 bL or bN.
@void@ WIR_UpDownValue::setBit@( unsigned int@ p, WIR_L4 v, const WIR_Location æ&æl @)@;Ü\label{lst:updnvalues4}Ü

// Get the L4 value of a single bit at a given position.
WIR_L4 æ&æWIR_UpDownValue::operator [] @( unsigned int@ p @)@;Ü\label{lst:updnvalues5}Ü

// Get the L4 value of a single bit at a given position.
const WIR_L4 æ&æWIR_UpDownValue::at@( unsigned int@ p @)@ const;Ü\label{lst:updnvalues6}Ü

// Determine whether an up/down value contains a given L4 value.
@bool@ WIR_UpDownValue::containsBit@(@ const WIR_L4 v @)@ const;Ü\label{lst:updnvalues7}Ü

// Determine whether an up/down value contains at least one L4 value from the given set.
@bool@ WIR_UpDownValue::containsBits@(@ WIR_L4Set æ&&æv @)@ const;Ü\label{lst:updnvalues8}Ü

// Determine whether an up/down value contains only the given L4 value.
@bool@ WIR_UpDownValue::containsOnlyBit@(@ const WIR_L4 v @)@ const;Ü\label{lst:updnvalues9}Ü

// Determine whether an up/down value contains only L4 values from the given set.
@bool@ WIR_UpDownValue::containsOnlyBits@(@ WIR_L4Set æ&&æv @)@ const;Ü\label{lst:updnvalues10}Ü

// Determine whether a bit at a given position is either bL or bN.
@bool@ WIR_UpDownValue::isLocationBit@( unsigned int@ p @)@ const;Ü\label{lst:updnvalues11}Ü

// Get the location at a given position.
const WIR_Location æ&æWIR_UpDownValue::getLocation@( unsigned int@ p @)@ const;Ü\label{lst:updnvalues12}Ü

// Get an up/down value's bit width.
@unsigned int@ WIR_UpDownValue::getBitWidth@( void )@ const;Ü\label{lst:updnvalues13}Ü

// Set a value's signedness.
@void@ WIR_UpDownValue::setSignedness@( bool@ s æ=æ true @)@;Ü\label{lst:updnvalues14}Ü

// Determine whether an up/down value is signed.
@bool@ WIR_UpDownValue::isSigned@( void )@ const;Ü\label{lst:updnvalues15}Ü

// Determine whether an up/down value is unsigned.
@bool@ WIR_UpDownValue::isUnsigned@( void )@ const;Ü\label{lst:updnvalues16}Ü
\end{lstlisting}

\cref{lst:updnvalues} shows all methods to set and retrieve the basic properties of up/down values, beginning with setting bits to some \texttt{L4} value.
Methods \texttt{setAllBits} (cf. \cref{lst:updnvalues1,lst:updnvalues2}) set all bits to the same given \texttt{L4} value, where the first variant of this method is applicable only for values other than \texttt{bL} and \texttt{bN}.
Since \texttt{bL} and \texttt{bN} necessarily require the specification of a location, the second variant is provided for this purpose.
Both methods produce a runtime error if wrong \texttt{L4} values are provided.
Similarly, methods \texttt{setBit} (\cref{lst:updnvalues3,lst:updnvalues4}) can be used to set a single bit specificed by its numerical position.
Bit positions within up/down values always start with 0 denoting the value's least-significant bit.
Retrieving the \texttt{L4} value of a bit at a given position can be done using the \texttt{[]} operator from \cref{lst:updnvalues5} or the \texttt{at} method (\cref{lst:updnvalues6}) if the up/down value is a \texttt{const} object.
Checking whether a given \texttt{L4} value occurs in some bit is done by method \texttt{containsBit} (\cref{lst:updnvalues7}).
By providing a set of \texttt{L4} bits to \texttt{containsBits} (\cref{lst:updnvalues8}), the existence of at least one of these values in some bits can be tested.
The set type \texttt{WIR\_L4Set} follows the principles discussed earlier in \cref{chap:setsmaps}.
Whether an up/down value contains only a given (set of) \texttt{L4} value(s) can be determined using \texttt{containsOnlyBit} (\cref{lst:updnvalues9}) and \texttt{containsOnlyBits} (\cref{lst:updnvalues10}), resp.
A convenient test whether a bit at a given position is a location \texttt{bL} or a negated location \texttt{bN} is realized by \texttt{isLocationBit} (\cref{lst:updnvalues11}).
The location at a bit position can be obtained using \texttt{getLocation} (\cref{lst:updnvalues12}).
It fails with a runtime error if the specified bit position is neither \texttt{bL} nor \texttt{bN}.

An up/down value's bit width can only be set during its construction.
For this reason, there is no method to change the bit width afterwards.
A value's bit width can be obtained using \texttt{getBitWidth} (\cref{lst:updnvalues13}).
The interpretation of an up/down value depends on whether it represents signed or unsigned contents.
A value's signedness can explicitly be set using \texttt{setSignedness} (\cref{lst:updnvalues14}) that expects a Boolean argument defaulting to \texttt{true}.
Methods \texttt{isSigned} and \texttt{isUnsigned} (\cref{lst:updnvalues15,lst:updnvalues16}) determine whether a value is signed or unsigned, resp.

\begin{lstlisting}[style=code,float=!t,caption={Value Interpretation by Class \texttt{WIR\_UpDownValue}},label=lst:updninterpret]
// Determine whether an up/down value contains only b0, b1 or bX bits.
@bool@ WIR_UpDownValue::isInteger@( void )@ const;Ü\label{lst:updninterpret1}Ü

// Determine whether an up/down value contains only b0 or b1 bits.
@bool@ WIR_UpDownValue::isBinaryInteger@( void )@ const;Ü\label{lst:updninterpret2}Ü

// Determine whether an up/down value is a positive integer number.
@bool@ WIR_UpDownValue::isPositive@( void )@ const;Ü\label{lst:updninterpret3}Ü

// Determine whether an up/down value is a negative integer number.
@bool@ WIR_UpDownValue::isNegative@( void )@ const;Ü\label{lst:updninterpret4}Ü

// Get an up/down value's numerical interpretation.
@signed long long@ WIR_UpDownValue::getSignedValue@( void )@ const;Ü\label{lst:updninterpret5}Ü
\end{lstlisting}

An up/down value can be interpreted as a classical number according to some number representation known from computer arithmetics.
However, the fact that \texttt{L4} values form a superset of standard Boolean bits needs to be considered.
\cref{lst:updninterpret} shows all methods of class \texttt{WIR\_UpDownValue} that can be used to numerically interpret an up/down value.
Method \texttt{isInteger} (cf. \cref{lst:updninterpret1}) determines whether an up/down value represents an integer number.
Integer numbers of course contain only bits \texttt{b0} and \texttt{b1}.
Since, however, \texttt{bX} bits are irrelevant and do not care by definition, such bits are also accepted by \texttt{isInteger}.
If a test for a standard binary integer shall be performed, method \texttt{isBinaryInteger} (\cref{lst:updninterpret2}) should be used that exclusively checks for \texttt{b0} and \texttt{b1} bits.
\texttt{isPositive} (\cref{lst:updninterpret3}) determines whether a value represents a positive number.
This is the case if the value either is unsigned, or if it is signed and it is an integer number whose most-significant bit (i.e., the sign bit) is \texttt{b0}.
Likewise, \texttt{isNegative} (\cref{lst:updninterpret4}) checks for a negative number, i.e., if a value is signed and is an integer with most-significant bit equal to \texttt{b1}.
An up/down value's numerical interpretation can be retrieved via method \texttt{getSignedValue} (\cref{lst:updninterpret5}).
The use of this method of course makes only sense if \texttt{isBinaryInteger} is \texttt{true}.
If an up/down value contains bits other than \texttt{b0} and \texttt{b1}, such bits are silently ignored, i.e., assumed to be \texttt{b0}.
For positive numbers, \texttt{getSignedValue} applies the standard unsigned number representation where all \texttt{b1} bits are multiplied by the power of 2 according to their bit positions, and all these partial products get accumulated.
For negative numbers, the two's complement representation is applied.
The so-determined numerical interpretation is then returned by this method as value of type \texttt{signed long long}, irrespective of the up/down value's actual signedness (cf. methods \texttt{isSigned} and \texttt{isUnsigned} from \cref{lst:updnvalues}).

\begin{lstlisting}[style=code,float=!t,caption={Comparison Operators of Class \texttt{WIR\_UpDownValue}},label=lst:updncompare]
// Determine whether two up/down values are provably, strictly bitwise equal.
WIR_L4 WIR_UpDownValue::operator == @(@ const WIR_UpDownValue æ&æv @)@ const;Ü\label{lst:updncompare1}Ü

// Determine whether two up/down values are provably, strictly bitwise inequal.
WIR_L4 WIR_UpDownValue::operator != @(@ const WIR_UpDownValue æ&æv @)@ const;Ü\label{lst:updncompare2}Ü

// Determine whether one up/down value is provably greater-equal than another.
WIR_L4 WIR_UpDownValue::operator >= @(@ const WIR_UpDownValue æ&æv @)@ const;Ü\label{lst:updncompare3}Ü

// Determine whether one up/down value is provably less than another.
WIR_L4 WIR_UpDownValue::operator < @(@ const WIR_UpDownValue æ&æv @)@ const;Ü\label{lst:updncompare4}Ü
\end{lstlisting}

Comparing up/down values for equality and other properties can be done using the operators shown in \cref{lst:updncompare}.
Again, the particularities of the \texttt{L4} partial order have to be considered so that these comparison operators behave differently from standard operators.
Due to the nature of \texttt{L4}, there can be situations where no clear result can be determined for a comparison.
For this reason, all the operators from \cref{lst:updncompare} do not return a simple Boolean result, but instead return an \texttt{L4} value that is either \texttt{b1} if a comparison definitely is true, or \texttt{b0} if a comparison provably is false, or \texttt{bU} if no safe statement can be made.

For the equality test of the \texttt{==} operator from \cref{lst:updncompare1}, the following rules apply:
\begin{itemize}
  \item If the compared up/down values have different bit widths, they are not equal and \texttt{b0} is returned.
  \item Otherwise, both values are equal and \texttt{b1} is returned iff, for each bit position, both compared bits are \texttt{b0} or both are \texttt{b1} or both are \texttt{bL} or both are \texttt{bN}, where, in the last two cases, both bits refer to exactly the same bit location.
  \item Otherwise, both values are inequal and \texttt{b0} is returned iff, for some bit position, one bit is \texttt{b0} and the other is \texttt{b1}, or one bit is \texttt{bL} and the other is \texttt{bN} where, in this case, both bits refer to exactly the same bit location.
  \item Otherwise, no clear statement can be made and \texttt{bU} is returned.
\end{itemize}
The analog rules (with inverted return values) are applied by the inequality operator \texttt{!=} from \cref{lst:updncompare2}.

The \texttt{>=} and \texttt{<} operators from \cref{lst:updncompare3,lst:updncompare4} are unable to make a safe statement and return \texttt{bU} if some bit of the two values to be compared is \texttt{bU}, \texttt{bL} or \texttt{bN}.
In other words, if some up/down value is not an integer (see \texttt{isInteger} from \cref{lst:updninterpret1} of \cref{lst:updninterpret}), both \texttt{>=} and \texttt{<} return \texttt{bU}.
If both values to be compared are integers, their numerical interpretations are obtained using \texttt{getSignedValue} and are compared with each other.
This way, either \texttt{b1} or \texttt{b0} are returned, depending on whether this numerical comparison holds true or not.

\begin{lstlisting}[style=code,float=!t,caption={Bitwise Operators of Class \texttt{WIR\_UpDownValue}},label=lst:updnbit]
// Bitwise AND operator.
WIR_UpDownValue WIR_UpDownValue::operator & @(@ const WIR_UpDownValue æ&æv æ)æ const;Ü\label{lst:updnbit1}Ü

// Bitwise OR operator.
WIR_UpDownValue WIR_UpDownValue::operator | @(@ const WIR_UpDownValue æ&æv æ)æ const;Ü\label{lst:updnbit2}Ü

// Bitwise XOR operator.
WIR_UpDownValue WIR_UpDownValue::operator Ü\^\ Ü@(@ const WIR_UpDownValue æ&æv æ)æ const;Ü\label{lst:updnbit3}Ü

// Bitwise negation operator.
WIR_UpDownValue WIR_UpDownValue::operator ~ @(@ const WIR_UpDownValue æ&æv æ)æ const;Ü\label{lst:updnbit4}Ü

// Bitwise left-shift operator by integer shift amount.
WIR_UpDownValue WIR_UpDownValue::operator << @( int@ s æ)æ const;Ü\label{lst:updnbit5}Ü

// Bitwise right-shift operator by integer shift amount.
WIR_UpDownValue WIR_UpDownValue::operator >> @( int@ s æ)æ const;Ü\label{lst:updnbit6}Ü

// Bitwise left-shift operator by another up/down value.
WIR_UpDownValue WIR_UpDownValue::operator << @(@ const WIR_UpDownValue æ&æv æ)æ const;Ü\label{lst:updnbit7}Ü

// Bitwise right-shift operator by another up/down value.
WIR_UpDownValue WIR_UpDownValue::operator >> @(@ const WIR_UpDownValue æ&æv æ)æ const;Ü\label{lst:updnbit8}Ü
\end{lstlisting}

Class \texttt{WIR\_UpDownValue} comes with the full set of bitwise and arithmetical operators in order to conveniently apply the usual computations to such up/down values and their featured \texttt{L4} bits.
\cref{lst:updnbit} shows the declarations of all bitwise operators.
Since most of them are straightforward, we omit their detailed discussion for the sake of brevity.
However, a few explanations regarding the shift operators from \crefrange{lst:updnbit5}{lst:updnbit8} are required.
The \texttt{<{}<} operator from \cref{lst:updnbit5} shifts an up/down value by the number of bits specified as integer shift amount.
If this shift amount is positive, a left-shift is applied, while negative shift amounts lead to a right-shift.
When shifting left, the \texttt{<{}<} operator fills the least-significant bits with \texttt{b0}.
When doing a right shift, the most-significant bits get filled with either \texttt{b0} (logical right-shift) or with the original up/down value's sign bit (arithmetical right-shift), depending on the value's signedness.
The same applies for the \texttt{>{}>} operator from \cref{lst:updnbit6} -- for positive shift amounts, either logical or arithmetical right-shifts are applied, while negative shift amounts are equivalent to a left-shift.

The shift operators from \cref{lst:updnbit7,lst:updnbit8} allow to use a shift amount being an up/down value itself.
If such a shift amount contains any \texttt{bU}, \texttt{bL} or \texttt{bN} bits, these shift operators produce a result containing only \texttt{bU} bits.
If, otherwise, the shift amount is positive, potential \texttt{bX} bits in the shift amount are interpreted as \texttt{b1}, or as \texttt{b0} for negative shift amounts.
This interpretation of \texttt{bX} bits maximizes the absolute value of the shift amount so that a maximal number of bits gets shifted out.
As before, left-shifting by a negative value de facto is a right-shift (either arithmetical or logical), and right-shifting by a negative value is a left-shift.

\begin{lstlisting}[style=code,float=ht,caption={Arithmetical and Miscellaneous Operators of Class \texttt{WIR\_UpDownValue}},label=lst:updnarith]
// Bitwise plus operator.
WIR_UpDownValue WIR_UpDownValue::operator + @(@ const WIR_UpDownValue æ&æv æ)æ const;Ü\label{lst:updnarith1}Ü

// Bitwise minus operator.
WIR_UpDownValue WIR_UpDownValue::operator - @(@ const WIR_UpDownValue æ&æv æ)æ const;Ü\label{lst:updnarith2}Ü

// Bitwise unary minus operator.
WIR_UpDownValue WIR_UpDownValue::operator - @( void )@ const;Ü\label{lst:updnarith3}Ü

// Bitwise multiplication operator.
WIR_UpDownValue WIR_UpDownValue::operator * @(@ const WIR_UpDownValue æ&æv æ)æ const;Ü\label{lst:updnarith4}Ü

// Bitwise division operator.
WIR_UpDownValue WIR_UpDownValue::operator / @(@ const WIR_UpDownValue æ&æv æ)æ const;Ü\label{lst:updnarith5}Ü

// Bitwise remainder operator.
WIR_UpDownValue WIR_UpDownValue::operator % @(@ const WIR_UpDownValue æ&æv æ)æ const;Ü\label{lst:updnarith6}Ü

// Get the absolute value.
friend WIR_UpDownValue abs@(@ const WIR_UpDownValue æ&æv @)@;Ü\label{lst:updnarith7}Ü

// Get the width-extended value.
WIR_UpDownValue WIR_UpDownValue::extend@( unsigned int@ w @)@ const;Ü\label{lst:updnarith8}Ü

// Get a consecutive bit sequence from an up/down value.
WIR_UpDownValue WIR_UpDownValue::extract@( unsigned int@ p, @unsigned int@ w @)@ const;Ü\label{lst:updnarith9}Ü

// Insert a value somewhere into another up/down value.
friend @void@ insert@(@ WIR_UpDownValue æ&æv1, const WIR_UpDownValue æ&æv2, @unsigned int@ p @)@;Ü\label{lst:updnarith10}Ü
\end{lstlisting}

Finally, \cref{lst:updnarith} summarizes all remaining operators of class \texttt{WIR\_UpDownValue} which mostly perform the standard arithmetical computations such as addition or subtraction.
The addition from \cref{lst:updnarith1} is realized as ripple-carry adder, while the subtraction (\cref{lst:updnarith2}) exploits the two's complement and thus gets reduced to an addition.
The unary minus operator (\cref{lst:updnarith3}) negates an up/down value by again applying the two's complement.
Multiplication (\cref{lst:updnarith4}) uses Booth's algorithm, while division and remainder from \cref{lst:updnarith5,lst:updnarith6} go back to a restoring division.

The \texttt{abs} function (\cref{lst:updnarith7}) provides the negated value for negative arguments, and the identity for positive arguments.
Extending the bit width of an up/down value can be achieved using \texttt{extend} (\cref{lst:updnarith8}).
This method gets the new bit width as argument (which must be larger than the value's current width) and returns either the zero- or the sign-extended up/down value.
From a bit position \texttt{p} on, \texttt{extract} (\cref{lst:updnarith9}) extracts the next \texttt{w} consecutive bits and returns this extracted value.
The \texttt{insert} function (\cref{lst:updnarith10}) takes the bits of value \texttt{v2} and inserts them into \texttt{v1} at bit position \texttt{p}.

The bit-true data flow analysis of class \texttt{WIR\_BitDFA} constructs the \gls{DFG} of a \gls{WIR} function according to \cref{def:bitdfg} and traverses it in forward and backward direction.
The first pass produces the down values $\mathit{Dn}$, while the second pass yields the $\mathit{Up}$ values.
Without going into the details of the analysis itself, the down values can never contain \texttt{bX} bits by construction.
Thus, the up values represent the final result of the whole analysis.
Besides the standard interface of \gls{WIR} analyses (cf. \cref{lst:wiranalysis}), this class does not provide any additional relevant methods.

Analysis results are stored in \texttt{WIR\_BitValues} containers (see file \texttt{WIR/containers/wirbitvalues.h}) which are attached to register parameters and to immediate parameters.
This is motivated by the fact that \gls{DFG} edges connect these types of parameters with each other as can be seen from \cref{fig:bitdfg}.\label{sec:bitinfo}

\begin{lstlisting}[style=code,float=ht,caption={Handling of Incoming Up/Down Values by Container Class \texttt{WIR\_BitValues}},label=lst:bitin]
// Add a down and an up value as incoming DFG edge.
@void@ WIR_BitValues::insertInValues@(@ WIR_Parameter æ&æp, WIR_UpDownValue æ&&æd, WIR_UpDownValue æ&&æu @)@;Ü\label{lst:bitin1}Ü

// Get the list of incoming DFG edges with their down and up values.
ästd::listäæ<æedgeæ> &æWIR_BitValues::getInValues@( void )@;Ü\label{lst:bitin2}Ü

// Find the incoming DFG edge starting at a given parameter.
ästd::listäæ<æedgeæ>::æiterator WIR_BitValues::findInValues@(@ const WIR_Parameter æ&æp @)@;Ü\label{lst:bitin3}Ü

// Remove the incoming DFG edge starting at the given parameter and its down and up values.
@void@ WIR_BitValues::eraseInValues@(@ const WIR_Parameter æ&æp @)@;Ü\label{lst:bitin4}Ü

// Remove the incoming DFG edge at the given iterator position and its down and up values.
@void@ WIR_BitValues::eraseInValues@(@ ästd::listäæ<æedgeæ>::æiterator pos @)@;Ü\label{lst:bitin5}Ü

// Remove all incoming DFG edges and their down and up values.
@void@ WIR_BitValues::clearInValues@( void )@;Ü\label{lst:bitin6}Ü

// struct edge represents a DFG edge between two WIR parameters.
struct WIR_BitValues::edgeÜ\label{lst:bitin7}Ü
{
  // rp points to the parameter being an edge's source or target node, resp.
  WIR_Parameter æ*ærp;Ü\label{lst:bitin8}Ü

  // downVal stores an edge's down value.
  WIR_UpDownValue downVal;Ü\label{lst:bitin9}Ü

  // upVal stores an edge's up value.
  WIR_UpDownValue upVal;Ü\label{lst:bitin10}Ü
};
\end{lstlisting}

Consider, e.g., register parameter \texttt{x\_3} of the \texttt{andi} operation from \cref{fig:bitdfg}.
This parameter is used by its operation so that data is flowing into it along an incoming \gls{DFG} edge.
The handling of incoming edges by container class \texttt{WIR\_BitValues} is shown in \cref{lst:bitin}.
For parameter \texttt{x\_3} of \texttt{andi}, the incoming edge gets added by the analysis using \texttt{insertInValues} (\cref{lst:bitin1}) where the incoming edge's source parameter (here: \texttt{x\_3} of \texttt{srli}) and the associated down and up values are provided.
This triplet of an incoming edge's source parameter and its up/down values is represented by struct \texttt{edge} from \cref{lst:bitin7}.
A reference to the complete list of incoming edges can be retrieved via \texttt{getInValues} (\cref{lst:bitin2}), while \texttt{findInValues} (\cref{lst:bitin3}) provides an iterator into this list referring to the edge having the given parameter as source.
Methods \texttt{eraseInValues} (\cref{lst:bitin4,lst:bitin5}) remove the incoming edge either starting at a given parameter or referred by a given iterator.
Finally, \texttt{clearInValues} (\cref{lst:bitin6}) erases all incoming edges.

The \texttt{andi} operation of the previous example from \cref{fig:bitdfg} produces register parameter \texttt{x\_4} as a result so that this register parameter is source of an outgoing \gls{DFG} edge.
In order to manage outgoing edges, container class \texttt{WIR\_BitValues} features the equivalent set of methods as already shown in \cref{lst:bitin}, which are just named \texttt{insertOutValues}, \texttt{getOutValues}, \texttt{findOutValues}, \texttt{eraseOutValues} and \texttt{clearOutValues}, resp.
Due to their similarity with \cref{lst:bitin}, we omit their detailed presentation here.
For an outgoing edge, struct \texttt{edge} from \cref{lst:bitin7} stores the register parameter being the edge's target node.


\cleardoubleoddpage
